% =====================================================================
%  capitulos/cap04.tex — Capítulo 4. La estructura interna de la palabra
%  Autor: Emiliano De Bin
%  Fuente: Capitulo_04_normalizado.docx
% =====================================================================
\capitulo{La estructura interna de la palabra}{Emiliano De Bin}

\section{Introducción}

En este capítulo presentamos y ejemplificamos los conceptos más relevantes de la
estructura interna de la palabra, de los que se ocupa la parte de la
lingüística que se conoce como morfología. Partimos del reconocimiento de la
existencia de unidades formales segmentables en el interior de las palabras,
establecemos luego las principales clases en las que se las agrupa y mostramos
cómo se organizan jerárquicamente estas unidades formando estructuras que
presentan relaciones de dependencia dentro de las piezas léxicas. Finalmente,
nos detenemos en los distintos procesos morfológicos, tanto los que intervienen
en la formación de palabras como los encargados de la manifestación de las
categorías gramaticales (género, número, tiempo, etc.).

\section{La segmentación morfológica}\label{sec:04-segmentacion}

En una primera aproximación intuitiva, a los hablantes del español nos resulta
evidente que la mayoría de las palabras de nuestra lengua no son un bloque
indivisible, sino que, por el contrario, son susceptibles de ser descompuestas
en unidades menores. Así, por ejemplo, en \forma{submarinas} podemos reconocer
cinco componentes: \forma{sub-}, \forma{-mar-}, \forma{-in-}, \forma{-a-} y
\forma{-s}, cada uno de los cuales contribuye con alguna parte de la
significación del ítem léxico en su conjunto. Estas unidades, que se denominan
\termino{morfemas}, se reconocen porque aparecen recurrentemente en otros ítems
léxicos de la lengua, realizando, por lo general, un aporte semántico semejante.
Encontramos a \forma{sub-} en \forma{subterráneo}, que comparte con
\forma{submarinas} la noción de \sig{estar debajo de}; \forma{mar}, por su
parte, puede conformar por sí solo una palabra, con el significado de
\sig{extensión de agua salada}; \forma{-in-}\footnote{Este segmento
\textit{-in-} es un sufijo adjetivador que aparece en palabras como
\textit{can-in-o} y \textit{cervantin-o}. No debe confundirse con el prefijo
\textit{in-} con valor negativo que se une a adjetivos: \textit{útil} >
\textit{inútil}.} reaparece en \forma{salino} con el valor de \sig{relativo a};
\forma{-a-} y \forma{-s}, como en \forma{alumna} y \forma{autos}, están
asociados a las ideas de \sig{género femenino} y \sig{pluralidad},
respectivamente. De modo análogo, \forma{descolorido} es susceptible de ser
descompuesto en \forma{des-}, que porta la idea de \sig{falta o privación} y
está presente en \forma{desnivel}; \forma{color} que, al igual que \forma{mar},
puede funcionar como una palabra independiente, referida a una propiedad física
de la materia; \forma{-id-}, que indica \sig{cualidad}, y reaparece en
\forma{dolorido}; y \forma{-o-}, vinculado a la noción de \sig{género
masculino}, como en \forma{alumno}.

La posibilidad de realizar esta descomposición habilita algunas preguntas. ¿Qué
realidad poseen los elementos que resultan de ella? ¿Se encuentran todos en el
mismo nivel jerárquico dentro de la palabra? ¿Juegan el mismo rol en su
funcionamiento interno? ¿De qué manera contribuyen estos formantes a la
predictibilidad del significado de los elementos léxicos de la lengua?

\subsection{Las dos ramas de la morfología: léxica y flexiva}\label{sec:04-ramas}

Interrogantes como los anteriores y otros que iremos presentando en este
capítulo constituyen el centro de interés de la morfología. Esta disciplina
gramatical, que estudia el funcionamiento interno de la palabra, está organizada
en dos ramas: la morfología léxica y la morfología flexiva.

Por un lado, la morfología léxica se pregunta por la estructura interna de la
palabra y por los procesos que permiten formar elementos léxicos nuevos sobre
la base de otros ya existentes. Por tanto, los procesos que permiten obtener
\forma{mar-in-o} o \forma{baja-mar} a partir de \forma{mar} son objeto de la
morfología léxica. La primera, \forma{mar-in-o}, se obtiene incorporando a la
forma simple de la palabra un sufijo, que también forma \forma{sal-in-o} a
partir de \forma{sal} o \forma{can-in-o} a partir de \forma{can}. Este proceso
se conoce como derivación. La segunda forma, \forma{bajamar}, proviene de la
unión de \forma{mar} con \forma{baja}, como también ocurre en
\forma{ple(n)a-mar} o \forma{alta-mar}. Se trata de un fenómeno de composición.
El español dispone de una gran variedad de procedimientos derivativos y
compositivos, como también de otros procesos menores de formación de palabras,
que iremos viendo a lo largo de este capítulo.

La morfología léxica también debe dar cuenta de cómo se encuentran
estructuradas por dentro las palabras. Por ejemplo, en el caso de
\forma{descolorido} está claro que el punto de partida para formarla es el
sustantivo \forma{color}, al que se une primero el sufijo \forma{-id-} para
formar \forma{colorido}, \sig{que tiene color}, y, luego, el prefijo
\forma{des-}, que construye el opuesto \forma{descolorido} (\sig{falto de
color}). Según vemos, en cada paso, la formación de la estructura se apoya en
una palabra existente. Sin embargo, hay casos un poco más complejos: así la
situación no es tan transparente en \forma{submarinas}. En el inicio de la
formación de esta palabra se encuentra el sustantivo \forma{mar}, pero, si bien
existe el adjetivo \forma{marino} (\sig{relativo al mar}) al cual podría
unírsele el prefijo \forma{sub-}, «tal análisis no es coherente con la
interpretación semántica del adjetivo:\footnote{En estos casos en que se produce
un desajuste entre la interpretación semántica y la segmentación morfológica,
se suele hablar de «paradojas de encorchetamiento». Dado el nivel introductorio
de esta obra, no profundizaremos en este tema, para lo que remitimos a
\textcite[4374-4376]{piera1999}.} \textit{submarino} no tiene el significado de
\sig{bajo lo relativo al mar} (interpretación que correspondería a la
yuxtaposición de \textit{sub-} y \textit{marino}: [\textit{sub} +
\textit{marino}]), sino \sig{relativo a lo que está bajo el mar} (interpretación
que implica que la estructura del adjetivo es [\textit{sub} + \textit{mar} +
\textit{ino}])» (Darmesteter, 1877, p.~129, cit. en
\cite[4734]{serranodolader1999}), que responde a un tipo de conformación
morfológica que veremos más adelante, en §\ref{sec:04-parasintesis}. Por tanto,
el desafío de la morfología léxica radica en poder explicar los pasos mediante
los que los morfemas se van relacionando entre sí, lo que, según vemos, puede
variar de una palabra a otra.

La morfología flexiva, por otro lado, da cuenta de las diversas formas que puede
adoptar una misma palabra para expresar distintas relaciones gramaticales y
adaptarse al funcionamiento sintáctico. Así, esta segunda vertiente explica, por
caso, que en el adjetivo \forma{submarinas} los morfemas \forma{-a-} y
\forma{-s} manifiestan dos categorías morfológicas que inciden en la combinatoria
sintáctica de la palabra para que pueda establecer relaciones con otras. Por
ejemplo, con un sustantivo masculino plural como \forma{recursos}, el adjetivo se
configura en una forma concordante para formar \forma{recursos submarinos}, la
cual replica el género y número del sustantivo al que modifica y, a la vez,
difiere de otra variante como \forma{submarina}, que es la utilizada para
concordar con un nombre femenino singular, como en \forma{fauna submarina}.
Además del género y el número, otras categorías que se expresan
morfológicamente en español son la persona, el tiempo, el modo y el aspecto.
Todas estas categorías no solo repercuten en la sintaxis: cada una de ellas
expresa algún contenido semántico muy general, como \sig{cantidad} o
\sig{tiempo}, tan relevante para la palabra que su omisión produce
agramaticalidad (v. \capref{2}): \juicio{*}{\forma{criaturas submarina}},
\juicio{*}{\forma{Llegará ayer}} (v. §\ref{sec:04-variantes}).

A lo largo de este capítulo abordaremos paulatinamente cada uno de estos
fenómenos. En todos los casos, el punto de partida del análisis morfológico es la
segmentación de la palabra en los componentes que la forman. A esta partición se
llega por medio de la comparación con otras palabras que exhiben los mismos
morfemas, esto es, observando su distribución dentro de las unidades léxicas.

\section{La unidad morfológica básica: el morfema}\label{sec:04-morfema}

Al igual que ocurre con otras unidades de la gramática, no resulta sencillo
definir con precisión qué es un morfema. En la tradición estructuralista, los
segmentos morfológicos fueron equiparados con el signo saussureano, tanto en la
vertiente norteamericana («forma lingüística con significado constante y
definido», \cite{bloomfield1933}) como en la europea, con la noción de monema
(«unidades dotadas de forma vocal y sentido», \cite{martinet1949}). Desde este
punto de vista, se consideró que era la asociación entre una representación
sonora y un contenido semántico lo que definía a los morfemas, los cuales
constituían «mínimas unidades significativas». Sin embargo, de la segmentación
morfológica de las palabras del español no siempre resultan unidades portadoras
de un contenido definido; así, por ejemplo, en los verbos, la llamada vocal
temática (VT) solo marca la clase conjugacional: \forma{-a} para los verbos de
primera conjugación (\forma{am-a-r}), \forma{-e} para los de segunda
(\forma{tem-e-r}), \forma{-i} para los de tercera (\forma{part-i-r}). Además,
algunos segmentos ni siquiera tienen una función clara en la estructura de la
palabra: se trata de componentes que funcionan como enlaces o que son resabios
de estados anteriores de la lengua. Por caso, en el análisis de una palabra como
\forma{inútil-i-dad-es}, el sufijo \forma{-dad} (\forma{mal-dad, bon-dad})
aparece precedido por una \forma{-i-} y, según veremos en el
ejercicio~\ref{ej:04-terraplanista}, en \forma{fals-e-dad-es}, el mismo sufijo
está precedido por una \forma{-e}, y en ambos casos la única función de estos
segmentos es vincular el sufijo \forma{-dad} con las respectivas bases
\forma{inútil-} y \forma{fals-}. Del mismo modo, también en el
ejercicio~\ref{ej:04-terraplanista}, la \forma{-a-} de \forma{plen-a-mente}
constituye un resto etimológico de la flexión de género que, en la actualidad,
no establece relación de concordancia ni recibe interpretación alguna.

Aceptado el hecho de que no todos los segmentos que se pueden obtener por medio
del análisis resultan significativos, otra objeción a la intuición que plantea a
los morfemas como unidades significativas es el hecho de que no siempre portan
un solo y único valor, sino que pueden aunar varios. Así, a veces un contenido
semántico se expresa de manera separada en un morfema independiente, como la
vocal temática en \forma{am-a-ban} y, en otros casos, ese contenido se presenta
unido en «amalgama» con otro segmento morfológico, como sucede con la vocal
temática en la primera persona singular del presente de indicativo, que se funde
con las otras categorías flexivas de la desinencia verbal (\forma{am-o, tem-o,
part-o}) (v.~§\ref{sec:04-vt}).

Por tanto, basándose en argumentos como los anteriores, \textcite[4318]{pena1999}
defiende una definición más amplia del morfema, como «la unidad mínima del
análisis morfológico y, en definitiva, del análisis gramatical», en el sentido
de que se trata de una unidad identificable dentro de la palabra, tenga o no
significado. En este último caso puede tener alguna de estas tres funciones:

\begin{itemize}
  \item distintiva o sin significado preciso (\forma{conocer / re-conocer}),
  \item clasificatoria (como la VT),
  \item formativa (como la \forma{-a} que aparece en algunos adverbios en
        \forma{-mente}).
\end{itemize}

\begin{ejercicio}[Terraplanistas]\label{ej:04-terraplanista}
Lea atentamente el siguiente fragmento extraído de una nota del diario español
\textit{El País} y luego realice las actividades indicadas más abajo.

% El fragmento del artículo se incorpora a mano desde el original (texto protegido).
\textoaincorporar{el fragmento de «No puedes convencer a un terraplanista y eso
debería preocuparte» (\textit{El País}, 27/02/2019,
\url{https://elpais.com/elpais/2019/02/27/ciencia/1551266455_220666.html}), con
las palabras subrayadas: \textit{convencer, terraplanista, negar, más,
falsedades, achatada, gigantesca, plenamente, representativos, explican,
entronización, pseudocientíficos}}

\textbf{Parte A}
\begin{consignas}
  \item Segmente en sus morfemas constitutivos las palabras subrayadas en el
        texto.
  \item Justifique el reconocimiento de cada morfema aportando al menos otras
        dos palabras en las que dicho morfema se encuentre presente con el mismo
        significado.
  \item Ofrezca una paráfrasis del significado y analice el aporte semántico que
        cada morfema realiza en los casos de \forma{falsedades},
        \forma{convencer}, \forma{plenamente} y \forma{explican}.
\end{consignas}

\begin{resuelto}
a) \forma{con-venc-e-r}. b) \forma{con-}: \forma{con-tribuir, con-tratar};
\forma{-venc-}: \forma{venc-iste, in-venc-ible}; \forma{-e-}: \forma{tem-e-rías,
respond-e-mos}; \forma{-r}:\footnote{El morfema \textit{-r} es el sufijo que
indica que se trata de una forma verbal de infinitivo.} \forma{ama-r,
parti-r}. c) \forma{fals-e-dad-es}. Significado: \sig{(dichos) que no se ajustan
a la verdad}. Aporte de cada morfema: \forma{fals-} \sig{que no se ajusta a la
verdad}; \forma{-dad-} \sig{propiedad de}; \forma{-es} \sig{plural}.
\end{resuelto}

\textbf{Parte B}

Decida cuáles de las generalizaciones expresadas más abajo describen con mayor
precisión lo observado.

\begin{consignas}
  \item Con respecto a la significatividad de los morfemas:
        \begin{exe}
          \ex Todos los morfemas que se pueden segmentar por medio del análisis
              distribucional realizan siempre un aporte semántico reconocible a
              la palabra.
          \ex No todos los morfemas que se pueden segmentar por medio del
              análisis distribucional realizan siempre un aporte semántico
              reconocible a la palabra: algunos de ellos tienen función de
              enlace entre otros morfemas o son vestigios de estadios anteriores
              de la lengua.
        \end{exe}
        \begin{resuelto}
        Opción correcta: 2.
        \end{resuelto}
  \item Con respecto a los límites del análisis formal:
        \begin{exe}
          \ex El análisis formal permite reconocer formantes morfológicos que se
              integraron a la palabra en otros estados de lengua anteriores y
              cuyo significado no es reconocible
              sincrónicamente.\footnote{\textit{Sincronía} y \textit{diacronía}
              son términos que provienen de \textcite{saussure1945}. La
              sincronía estudia un fenómeno tal como es en un momento
              específico, mientras que la diacronía estudia la evolución de ese
              mismo fenómeno a lo largo del tiempo.}
          \ex El análisis formal permite reconocer solamente formantes
              morfológicos integrados a la palabra sincrónicamente, y cuyo
              significado resulta siempre reconocible.
        \end{exe}
        Opción correcta:\completar[1]
  \item Con respecto a la relación entre los segmentos reconocidos y sus
        significados:
        \begin{exe}
          \ex Existen segmentos morfológicos que pueden expresar distintos
              significados y significados que pueden expresarse por medio de
              distintos segmentos morfológicos.
          \ex Todos los segmentos morfológicos expresan siempre los mismos
              significados, pero no todos los significados se expresan siempre
              en los mismos segmentos morfológicos.
          \ex Existen segmentos morfológicos que pueden expresar distintos
              significados, pero cada significado se expresa siempre en el mismo
              segmento morfológico.
          \ex Todos los segmentos morfológicos aportan siempre significados
              distintos.
        \end{exe}
        Opción correcta:\completar[1]
\end{consignas}
\end{ejercicio}

\subsection{La alomorfía}\label{sec:04-alomorfia}

Algunos morfemas del español están asociados a más de una manifestación
fonológica o morfo. Estas variantes de realización, que guardan entre sí un
parecido formal y se asocian todas de manera estable al mismo significado, se
denominan \termino{alomorfos}.

La alomorfía puede originarse en causas diversas. En algunos casos, la
alternancia puede estar motivada por factores externos al sistema gramatical.
Por ejemplo, por haber aparecido como resultado del cambio lingüístico de la
palabra. Así, las palabras \forma{apert-ura}, \forma{abert-ura} y
\forma{abr-i-r} tienen su origen etimológico en la raíz verbal latina
\forma{aper-}, que sufrió la sonorización de la consonante sorda intervocálica.
Eso produjo la coexistencia entre los alomorfos \forma{aper-}, \forma{aber-} y
\forma{abr-} (este último con caída de la segunda vocal). Dicho de otro modo,
estas variantes se explican, en última instancia, por razones léxicas.

Sin embargo, la motivación para la alomorfía también puede provenir del propio
sistema gramatical. En este sentido, hay morfemas cuya manifestación se ve
condicionada por factores vinculados con el contexto morfológico en el que
aparecen. Por ejemplo, por los morfemas que lo rodean: \forma{am-a-ble},
\forma{am-a-bil-ísim-o}, \forma{am-a-bil-idad} y \forma{am-a-ble-mente}
despliegan dos variantes en alternancia para un mismo sufijo: \forma{-ble-} y
\forma{-bil-}. La motivación para la aparición de una u otra forma radica en que
\forma{-ble} se emplea en posición final de palabra, mientras que \forma{-bil-}
aparece cuando le sigue(n) otro(s) morfema(s).

La alomorfía también puede venir motivada por razones fonológicas. Por caso, en
muchas palabras del español las vocales medias (\forma{e}, \forma{o}) diptongan
cuando reciben acento léxico. Es lo que ocurre en el contraste entre
\forma{hel-a-r} y \forma{hiel-o}. De la misma manera, el participio
\forma{abier-t-o} exhibe una cuarta forma de la raíz ejemplificada
anteriormente: \forma{abier-}. Lo que causa la aparición de esta última variante
es, en este caso, un factor de naturaleza fonológica que abre la vocal cuando
esta recibe el acento.

\begin{ejercicio}\label{ej:04-alomorfos}
Considere el siguiente corpus, que recoge palabras extraídas del texto del
comienzo del capítulo y lo completa con otras. Todas ellas manifiestan
alomorfos.

\begin{consignas}
  \item Identifique los morfemas que presentan alomorfos y describa su
        significado.
  \item Explique la causa que motiva la aparición de una u otra variante.
\end{consignas}

\begin{tabular}{@{}llll@{}}
  \forma{cuesta}        & \forma{terraplanistas} & \forma{culturales}  & \forma{irrumpido} \\
  \forma{terráqueos}    & \forma{roto}           & \forma{datos}       & \forma{costosa} \\
  \forma{terrícola}     & \forma{entierran}      & \forma{falsedades}  & \forma{ruptura} \\
  \forma{jóvenes}       & \forma{territorio}     & \forma{irrupción}   & \forma{disruptiva} \\
  \forma{maniqueas}     & \forma{aterrizar}      & \forma{rompió}      & \forma{expertos} \\
  \forma{Tierra}        & \forma{(las) crisis}   & \forma{costar}      & \forma{cuesten} \\
\end{tabular}

\begin{resuelto}
a) \forma{cuest-a, cost-a-r, cuest-e-n, cost-os-a}. En las palabras de este
grupo se identifica el morfema \textsc{cost-}, que significa \sig{precio o
esfuerzo demandado}. Los alomorfos son \fonol{ˈkwest} \textasciitilde{}
\fonol{kost}.\footnote{El signo \textasciitilde{} se lee \sig{alterna con}.} b)
La causa de la alomorfía es fonológica: \fonol{ˈkwest} aparece en sílaba tónica
y \fonol{kost}, en sílaba átona.
\end{resuelto}
\end{ejercicio}

\subsection{Asimilación y disimilación}

Aunque cada nivel lingüístico posee sus unidades y operaciones características,
no son pocos los fenómenos en los que estos entran en contacto. Algunos casos de
alomorfía fonológicamente motivada resultan de procesos fonológicos, como los de
asimilación y disimilación. El primero, que ya se describió en el \capref{3},
§\ref{sec:03-procesos}, consiste en que dos segmentos fónicos se vuelven
articulatoriamente más parecidos entre sí cuando entran en contacto. El caso más
típico es el de las consonantes nasales que, cuando preceden a otros sonidos
consonánticos, adaptan su punto de articulación al del segmento que les
sigue.\footnote{Es lo que se conoce como asimilación regresiva y es el tipo de
asimilación más frecuente en las lenguas (v. \capref{3},
§\ref{sec:03-procesos}).} De esto resulta que morfemas como el prefijo negativo
\forma{in-}, cuando se añaden a palabras iniciadas en consonante bilabial, se
realicen con una nasal bilabial: \forma{im-batible} \fon{imbaˈtiβle}, mientras
que en \forma{in-fum-a-ble} \fon{iɱfuˈmaβle}, \forma{in-toler-a-ble}
\fon{intoleˈɾaβle} e \forma{in-com-i-ble} \fon{iŋkoˈmiβle}, aunque la ortografía
no cambia e invisibiliza el fenómeno, también se pronuncian alófonos distintos
del archifonema nasal (v. \capref{3}, §\ref{sec:03-procesos}). El panorama de
posibilidades se completa ante palabras cuya primera consonante es un sonido
líquido: en este caso, el sonido nasal del prefijo cae y no se produce:
\forma{i-legal}, \forma{i-rracional}.\pendiente{El original dice «sonido
lateral», pero incluye \textit{irracional}}

Mientras que en la asimilación los sonidos se tornan formalmente más semejantes,
en la disimilación ocurre el proceso inverso: un segmento fonológico cambia su
forma habitual para diferenciarse de otro cercano. Las consonantes líquidas
suelen participar de este fenómeno. Eso explica por qué \forma{invern-al} y
\forma{circul-ar}, que contienen el mismo sufijo derivativo que expresa un valor
relacional (respectivamente \sig{relativo al invierno} y \sig{relativo al
círculo}), manifiestan, en el primer caso, un alomorfo terminado en lateral
(\forma{-al}) y, en el segundo, uno que finaliza en vibrante (\forma{-ar}) para
evitar la reiteración del segmento lateral: \juicio{*}{\forma{circulal}},
pronunciación que resultaría cacofónica en español.

\section{Clases de morfemas}

Los morfemas se pueden agrupar en clases según dos criterios principales: por
sus propiedades distribucionales y por su funcionamiento en el interior de la
palabra.

\subsection{Morfemas libres y ligados}

Desde un punto de vista distribucional, los morfemas se dividen en dos grandes
grupos: los \termino{morfemas libres}, que pueden aparecer solos ---o sea, tienen
capacidad para formar por sí solos una palabra--- y los \termino{morfemas
ligados}, que siempre aparecen en concurrencia con otros. Así, decimos que el
segmento \forma{-mar-}, que integra la palabra \forma{sub-mar-in-a-s}, pertenece
al primer grupo, dado que, como vimos, por sí mismo constituye una forma léxica
que puede funcionar dentro de una oración: \forma{Caminé hacia el mar}; en
cambio, ni el sufijo \forma{-in-} ni el prefijo \forma{sub-} ofrecen esta
posibilidad: su aparición siempre depende de que se encuentren junto con otros
segmentos morfológicos, como en \forma{bov-ino} o \forma{sub-desarrollado}.

\subsection{Raíces y afijos}

Más allá de su distribución, no todas las unidades morfológicas tienen la misma
función en la palabra. Las \termino{raíces}, que son las responsables del
significado léxico fundamental de la palabra, pueden ser libres, como
\forma{mar, árbol} o \forma{sol}, o ligadas, como \forma{convenc-} (en
\forma{convencer}) o \forma{chat-} (en \forma{achatado}). La mayor parte de las
raíces del español son ligadas y solo constituyen formas libres las que
integran palabras monomorfemáticas, como las de clase cerrada (categorías
funcionales, preposiciones, adverbios no terminados en \forma{-mente}) y un
puñado de adjetivos y nombres terminados en consonante, como \forma{árbol},
\forma{señor} o \forma{reloj}. Los miembros de las llamadas «familias de
palabras» comparten la raíz. Por ejemplo, \forma{submarino} se agrupa por medio
de este criterio con las ya mencionadas \forma{marino, bajamar, pleamar,
altamar} y también con \forma{marea, marear, marítimo, marinero, ultramarino,
marejada, maricultura, mareógrafo}, entre otras.

Los \termino{afijos} son unidades siempre ligadas entre las que cabe
diferenciar, por un lado, los derivativos, que aportan información léxica y
permiten formar palabras nuevas (como \forma{-in-} en \forma{mar-in-o} o
\forma{anti-} en \forma{anti-hongo-s}) y, por otro lado, los flexivos, que son
los que brindan información gramatical vinculada a las categorías lingüísticas
(como \forma{-o-} en \forma{marin-o}, que expresa género \sig{masculino}, y
\forma{-s} en \forma{antihongo-s}, que manifiesta \sig{plural}). Dependiendo de
su posición relativa respecto de la raíz, los afijos se dividen en prefijos y
sufijos. En español observamos que los flexivos se realizan exclusivamente como
sufijos, dado que nuestra lengua no posee prefijos
flexivos.\pendiente{El original dice «los derivativos se realizan exclusivamente
como sufijos»} En cuanto al orden relativo, en una palabra como
\forma{gigantesca} podemos observar que la raíz \forma{gigant-} primero se une
al sufijo derivativo \forma{-esc-} para formar el adjetivo \forma{gigantesc-}; y
luego entra en juego la morfología flexiva que incorpora el sufijo de género
femenino \forma{-a-}. Y, si quisiéramos pluralizar la palabra, el sufijo flexivo
plural \forma{-s} se agregaría después del de género. La secuencia que acabamos
de ejemplificar no es casual ni aislada: los sufijos flexivos son siempre más
externos que los derivativos en relación con la raíz, ya que, como demostró
\textcite{bybee1985}, los derivativos afectan el significado de la raíz,
mientras que los flexivos inciden en el comportamiento sintáctico de la palabra.
Asimismo, los elementos flexivos no solo son los últimos en incorporarse a la
palabra, sino que, además, lo hacen en un orden determinado. Por ejemplo, en una
forma verbal como \forma{estudi-a-ba-n}, a la raíz \forma{estudi-} le sigue la
vocal temática \forma{-a-}, a continuación se añade la manifestación amalgamada
de modo indicativo, tiempo pretérito y aspecto imperfectivo, manifestada por el
sufijo \forma{-ba-} y, finalmente, se une otra amalgama, que indica
conjuntamente tercera persona y número plural: \forma{-n}
(v.~§\ref{sec:04-variantes}).

Más allá de su evidente diferencia posicional, prefijos y sufijos poseen
características de funcionamiento distintas. Así, mientras que los primeros no
suelen ser selectivos con relación a la clase de la palabra a la que se
incorporan, los segundos son sensibles a la categoría léxica y, en general, la
modifican. Como ya hemos comentado, \forma{marino} se obtiene de \forma{mar} más
el sufijo \forma{-in}, que forma adjetivos relacionales a partir de nombres
(\forma{salino} < \forma{sal}, \forma{canino} < \forma{can}).\footnote{El signo <
se lee \sig{proviene de}.} Por contraste, el prefijo \forma{sub-} produce un
nombre a partir de otro en \forma{subgrupo} (< \forma{grupo}) y
\forma{subsistema} (< \forma{sistema}), pero también verbos derivados a partir
de verbos, como en \forma{subestimar} (< \forma{estimar}) y \forma{subarrendar}
(< \forma{arrendar}). Fonológicamente, el nivel de integración que exhiben con
la palabra también varía: suele ser mayor en el caso de los sufijos que en el de
los prefijos, entre los que existe incluso un grupo de formas separables, con
acento léxico propio, como \forma{ex-} en \forma{ex primer ministro} o
\forma{re-} en \forma{re me cae bien}.\footnote{Según la \textit{Ortografía de la
lengua española} \parencite[cap.~V, §\,2.2.2]{rae2010}, los prefijos deben
siempre escribirse unidos a la base, excepto, como en los casos de arriba,
cuando se coloquen delante de varias palabras.}

\begin{ejercicio}\label{ej:04-clases}
\textbf{Parte A}
\begin{consignas}
  \item Complete la siguiente tabla, clasificando los morfemas segmentados en el
        ejercicio~\ref{ej:04-terraplanista} a partir de su funcionamiento en la
        palabra.

        {\footnotesize
        \begin{tabular}{@{}l*{5}{c}@{}}
          \toprule
          Palabra & \makecell{Pref.\\flexivo} & \makecell{Pref.\\derivativo} & Raíz & \makecell{Suf.\\derivativo} & \makecell{Suf.\\flexivo} \\
          \midrule
          \forma{convencer}      & & con- & -venc- & -e- & -r \\
          \forma{falsedades}     & & & fals- & -e-dad- & -es \\
          \forma{terraplanistas} & & & & & \\
          \ldots & & & & & \\
          \bottomrule
        \end{tabular}}
  \item Identifique cuáles de los morfemas segmentados pueden clasificarse como
        libres y cuáles solo pueden funcionar como ligados.
\end{consignas}

\textbf{Parte B}

A partir de lo observado, decida cuáles de las siguientes generalizaciones
describen con mayor precisión para el español el funcionamiento de los morfemas
analizados.

\begin{consignas}
  \item Respecto de la relación entre las clases distribucionales y los tipos de
        morfemas:
        \begin{exe}
          \ex Las raíces son siempre morfemas libres, pero no todos los
              morfemas libres son siempre raíces.
          \ex Tanto las raíces como los afijos son siempre morfemas ligados.
          \ex Los afijos pueden ser morfemas libres o morfemas ligados.
          \ex Los afijos son siempre morfemas ligados, pero no todos los
              morfemas ligados son siempre afijos.
        \end{exe}
        \begin{resuelto}
        Opción correcta: 4.
        \end{resuelto}
  \item Respecto del aporte gramatical que realizan los sufijos y los prefijos
        del español a la construcción a la que se integran:
        \begin{exe}
          \ex Tanto los prefijos como los sufijos pueden expresar valores
              flexivos.
          \ex Los sufijos siempre forman palabras nuevas, pero no todas las
              palabras nuevas se forman por medio de la sufijación.
          \ex Solo los sufijos pueden expresar valores flexivos.
          \ex Solo los prefijos permiten formar palabras nuevas.
        \end{exe}
        Opción correcta:\completar[1]
\end{consignas}
\end{ejercicio}

\subsection{Otros afijos}

Existen otras formas de afijos que aparecen con menor frecuencia en la
morfología del español. Por un lado, hay formantes que se ubican entre raíces y
prefijos o entre raíces y sufijos (de los que ya hemos hablado en
§\ref{sec:04-morfema}), cuya función fundamental es actuar como enlaces
fonológicos, por lo que no portan significado. Es lo que ocurre con el segmento
\forma{-c-}, que se inmiscuye entre la raíz \forma{café-} y el sufijo
\forma{-it-} en \forma{cafe-c-ito}. Del mismo modo funciona el segmento
\forma{-t-} en \forma{café-t-ería} ante el sufijo \forma{-ería} y ante
\forma{-iv-} en \forma{representa-t-ivos}. A estos morfemas se los conoce como
\termino{interfijos}.

Por otro lado, existen casos muy marginales para el español en los que un
morfema puede partir la raíz e insertarse en su interior. A estos segmentos se
los denomina \termino{infijos}, y se limitan, en nuestra lengua, a la expresión
de algunos diminutivos, como en \forma{azuqu-ít-ar} y \forma{lej-it-os}.

Finalmente, existen afijos que constituyen morfemas discontinuos. Los
\termino{circunfijos}, por ejemplo, son formantes que se aplican rodeando la
palabra a la que se añaden. Puede pensarse que \forma{tras-papel-a-r} se
encuentra formado así, ya que ni \forma{tras-} ni \forma{-a-} se han incorporado
a la raíz \forma{papel} por separado, pues no se registran ni
\juicio{*}{\forma{traspapel}} ni \juicio{*}{\forma{papelar}}. Sin embargo, se
trata de formantes que funcionan de manera independiente en otras palabras,
como \forma{tras-nacion-al} o \forma{sal-a-r}, por lo que cabe mejor pensarlos
como prefijos y sufijos que se incorporan en un mismo paso en el proceso de
formación de la palabra. Analizaremos estos casos con más detalle más adelante,
al explicar el proceso de formación de palabras conocido como parasíntesis
(v.~§\ref{sec:04-parasintesis}).

\subsection{El concepto de \textit{tema}}\label{sec:04-tema}

Dentro de cada elemento léxico, la unidad a partir de la cual se forman las
variantes flexivas se denomina \termino{tema}. El tema es la unidad que codifica
el contenido léxico de la palabra y es común a todas sus formas flexionadas. Si
retomamos nuestro ejemplo de §\ref{sec:04-segmentacion}, el segmento
\forma{marin-} allí considerado constituye, precisamente, el tema de la palabra
\forma{marinos}, ya que es a \forma{marin-} que se unen los sufijos flexivos
\forma{-o-} de género y \forma{-s} de plural. Asimismo, el tema \forma{norma-}
es la unidad que permite formar tanto el singular \forma{norma} como el plural
\forma{norma-s}. A partir del tema \forma{norma-} se puede derivar
\forma{normal-} (\sig{relativo a la norma}) y, a partir del nuevo tema
\forma{normal-}, se producen tanto el plural \forma{normal-es} como el derivado
adverbial \forma{norm-al-mente}, mediante el agregado de un sufijo flexivo de
plural, en el primer caso, y de uno de derivación, en el segundo. Estos procesos
de derivación pueden representarse mediante una descripción estructural en la
que cada corchete encierra un paso en la formación de la palabra e indica la
clase de palabra resultante, como se muestra en \exref{ex:04-normalmente}:

\begin{exe}
  \ex\label{ex:04-normalmente} Descripción estructural (DE):
      [[[norma]\textsubscript{N} al]\textsubscript{A} mente]\textsubscript{Adv}\footnote{Las
      abreviaturas remiten a las clases de palabras resultantes de cada etapa:
      N: nombre; A: adjetivo; Adv: adverbio; V: verbo.}
\end{exe}

Entre los morfemas que integran el tema de la palabra existen relaciones
jerárquicas, lo que significa que expresan una dependencia de tipo estructural.
Esto implica, entre otras cosas, que en la forma de la palabra queda un
testimonio de los pasos que permitieron constituirla. Se considera que cada uno
de estos pasos debe conformar una palabra existente en la lengua (aunque más
adelante veremos que en algunos procesos este requisito no se mantiene,
v.~§\ref{sec:04-parasintesis}), que puede, a su vez, ser flexionada de acuerdo
con las posibilidades de su categoría.

Operativamente se puede decir que el tema resulta de quitarle a la palabra todos
sus afijos flexivos. En consecuencia, los afijos flexivos, que en español se
realizan exclusivamente como sufijos, no integran el tema. Esto no implica, no
obstante, que no pueda haber afijos derivativos en el tema, tanto sufijos como
prefijos. Por tanto, dado que la flexión es un proceso siempre posterior a la
derivación, cada paso en la historia derivacional de la palabra coincide con un
tema, y solo en la última etapa morfológica ---concluido ya el ciclo de
formación de la palabra--- se pueden adjuntar morfemas flexivos. Como se muestra
en la descripción estructural \exref{ex:04-marinos}, al sustantivo \forma{mar}
se le une el sufijo \forma{-in-} para formar un adjetivo y a este tema se le
añaden los elementos flexivos, que no involucran un paso más (otro corchete)
porque su agregado no involucra la formación de una palabra distinta.

\begin{exe}
  \ex\label{ex:04-marinos} DE: [[mar]\textsubscript{N} in]\textsubscript{A}-o-s
\end{exe}

\begin{ejercicio}\label{ej:04-tema}
\textbf{Parte A}

A partir de las palabras \forma{achatada} y \forma{pseudocientíficos}, que
fueron segmentadas en el ejercicio~\ref{ej:04-terraplanista}, e incorporando al
corpus \forma{entronización, culturales, bromear} y \forma{desprecio}, tomadas
del mismo texto, realice las actividades que se indican a continuación.

\begin{consignas}
  \item Segmente los ejemplos que no fueron analizados previamente en morfemas y
        proporcione una paráfrasis del significado de cada palabra.
  \item Indique los pasos mediante los que se ha ido formando la palabra y a qué
        categoría pertenece el tema resultante en cada uno de ellos.
  \item Realice la descripción estructural (DE) de la palabra.
\end{consignas}

\begin{resuelto}
\forma{entronización}: a) \forma{en-tron-iz-a-ción} = \sig{acción o efecto de
poner algo o a alguien en un trono}. b) \forma{trono}\textsubscript{N} >
\forma{entronizar}\textsubscript{V} > \forma{entronización}\textsubscript{N}; N >
V > N\textsubscript{abstracto}. c) [[en-[tron]\textsubscript{N}-iz-a]\textsubscript{V}-ción]\textsubscript{N}.
\end{resuelto}

\textbf{Parte B}

Decida cuáles de las siguientes afirmaciones describen con mayor precisión lo
observado en el análisis estructural realizado.

\begin{consignas}
  \item Con respecto al tema:
        \begin{exe}
          \ex Operativamente, el tema resulta de quitarle a la palabra todos sus
              afijos flexivos.
          \ex Operativamente, el tema resulta de quitarle a la palabra todos sus
              sufijos.
          \ex Operativamente, el tema resulta de quitarle a la palabra todos los
              afijos.
          \ex Operativamente, el tema coincide siempre con la raíz de la
              palabra.
        \end{exe}
        \begin{resuelto}
        Opción correcta: 1.
        \end{resuelto}
  \item Con respecto al proceso derivacional:
        \begin{exe}
          \ex Cada paso del proceso derivacional coincide con la incorporación
              de un morfema.
          \ex Cada paso del proceso derivacional coincide con la incorporación
              de un morfema derivativo.
          \ex Cada paso del proceso derivacional coincide con el tema de una
              palabra existente.
          \ex Cada paso del proceso derivacional coincide con una palabra
              flexionada.
        \end{exe}
        Opción correcta:\completar[1]
\end{consignas}
\end{ejercicio}

\subsection{Los temas cultos o semipalabras}\label{sec:04-temascultos}

Existe un segundo sentido según el cual se utiliza el término \forma{tema} en
español, que se diferencia del que hemos utilizado hasta ahora. Designa a una
clase de morfemas asociados a contenidos léxicos, los cuales por sí solos no
pueden recibir morfología flexiva ni funcionar como palabras, por lo que se los
conoce como \termino{semipalabras}, aunque también, por su origen en las lenguas
clásicas, se los llama \termino{temas cultos}. A diferencia de las palabras, que
siempre pertenecen a una clase determinada (N, V, A, etc.), los temas cultos no
se categorizan hasta que, para formar una palabra, se unen con otro elemento de
la lengua. Ese segundo formante puede ser un afijo derivativo (\forma{a-fonía},
\forma{fon-iatra}), otra palabra (\forma{fono-terapia}) u otro tema culto
(\forma{fono-logía}), como veremos en §\ref{sec:04-composicion}.

\begin{recuadro}[Para ampliar]
Aunque las gramáticas clásicas los asimilaban a los prefijos, los temas cultos
no siempre van al inicio de la palabra, como sí lo hacen los prefijos. En
realidad, los temas cultos no tienen posición fija en la palabra, en la que
pueden estar al principio o al final, como lo muestran pares como
\forma{fono-logía / telé-fono}, \forma{podo-logía / cefaló-podo},
\forma{logo-pedia / astró-logo}.
\end{recuadro}

\subsection{La vocal temática}\label{sec:04-vt}

Como ya adelantamos en §\ref{sec:04-morfema}, en algunas formas flexivas de los
verbos del español se advierte la presencia de un segmento que carece de
contenido léxico, pero se caracteriza por indicar la clase según la cual se
conjuga el verbo. Se trata de la \termino{vocal temática}, que aparece, por
ejemplo, precediendo a la flexión de infinitivo, donde contrastan las tres
conjugaciones: \forma{-a-} en \forma{amar}, \forma{-e-} en \forma{temer},
\forma{-i-} en \forma{partir}. Este morfema se manifiesta aislado o fusionado con
otros rasgos flexivos. Resulta segmentable en algunas formas, como la primera
persona del plural del pretérito imperfecto del indicativo
(\forma{am-á-ba-mos, tem-í-a-mos, part-í-a-mos});\footnote{Como puede apreciarse
en estos ejemplos, la vocal temática de la segunda y la tercera conjugación se
manifiesta por medio del mismo morfo.} en cambio, en otras, como la primera
persona del singular del presente del indicativo (\forma{am-o, tem-o, part-o}),
no tiene manifestación independiente. Estos morfemas, que contienen más de un
significado, se denominan \textit{portmanteau} (francés: \sig{maleta}) o
\termino{amalgama}. Una tercera posibilidad es que la vocal temática sea el
único elemento que acompaña al tema. Es un análisis posible para lo que ocurre
en los imperativos de segunda persona del singular, en los que las otras
categorías flexivas (tiempo, modo, aspecto, número y persona) no se manifiestan
(\forma{am-á, tem-é, part-í}).

Como se desprende de lo anterior, dado que señala la clase conjugacional del
verbo, la vocal temática permite predecir cómo se obtienen determinadas
variantes flexivas, como, por ejemplo, los pretéritos imperfectos, ya que la
clase conjugacional del verbo explica por qué no producimos formas como
\juicio{*}{\forma{amíamos}}, \juicio{*}{\forma{temábamos}} o
\juicio{*}{\forma{partábamos}}.

\begin{recuadro}[Para ampliar: la VT en los nombres]
El análisis anterior se puede extender al ámbito nominal, ya que en sustantivos
y adjetivos la formación del plural también se determina por la naturaleza de
la terminación del tema, según los siguientes casos:
\begin{itemize}
  \item Tema finalizado en consonante, salvo que esté seguido por un sufijo
        flexivo de género: el plural se forma con la variante \forma{-es}:
        \forma{ventanal-es} (y no \juicio{*}{\forma{ventanals}}).
  \item Tema finalizado en \fonol{s} (palabras graves y esdrújulas): las formas
        singulares y plurales son homónimas: \forma{la crisis / las crisis}.
  \item Tema finalizado en vocal o seguido por la marca de género: la flexión se
        realiza con \forma{-s} (\forma{casa-s, gat-a-s}, pero no
        \juicio{*}{\forma{casa-es}}, \juicio{*}{\forma{gat-a-es}}).
  \item Tema finalizado en vocales cerradas acentuadas: forma el plural como los
        temas finalizados en consonante, \forma{ñandú-es} y \forma{gurí-es}. No
        obstante, para estas terminaciones algunas variedades pueden admitir el
        alomorfo \forma{-s} (\forma{ñandú-s, gurí-s}) o incluso la variante
        \forma{-ses} (\forma{ñandú-ses, gurí-ses}).
\end{itemize}
\end{recuadro}

\section{Homonimia estructural en el nivel morfológico}\label{sec:04-homonimia}

Entre los morfemas que integran el tema existen relaciones de dependencia, por
lo que también algunas variantes interpretativas de las palabras pueden
explicarse como casos de homonimia estructural. Esto quiere decir que su
estructura se puede describir por medio de (al menos) dos análisis
alternativos. Por ejemplo, la palabra \forma{in-mov-il-iz-a-ble} puede
significar \sig{que se puede inmovilizar} o \sig{que no se puede movilizar}. La
interpretación que realizamos depende de en qué nivel de la formación de la
palabra consideramos que se adjunta el prefijo, si al tema \forma{móvil} o a
\forma{movilizable}, lo cual, según vemos, no es trivial, ya que cambia el
significado con que interpretamos la palabra. En la primera estructura, el tema
del verbo \forma{inmovilizar} se forma a partir del adjetivo \forma{inmóvil} y la
incorporación del sufijo \forma{-ble} produce el significado \sig{que se puede
inmovilizar}. En la segunda interpretación, \forma{in-} se une al ya formado
adjetivo \forma{movilizable} y genera la lectura \sig{no movilizable}, es decir,
\sig{que no se puede movilizar}. En definitiva, los segmentos morfológicos que
intervienen son los mismos, pero la trama de relaciones entre ellos varía y, en
consecuencia, también lo hace la interpretación.

\begin{ejercicio}\label{ej:04-homonimia}
Partiendo de que las siguientes palabras tienen por lo menos dos acepciones cuyo
origen puede explicarse estructuralmente, realice las actividades que se
solicitan a continuación:

\begin{consignas}
  \item Deduzca a partir de los ejemplos presentados las dos interpretaciones
        alternativas asociadas a cada una de estas palabras y realice una
        paráfrasis de cada uno de los significados propuestos.
  \item Presente los pasos de la derivación y ofrezca descripciones
        estructurales que den cuenta de los significados asignados a cada
        palabra.
\end{consignas}

\begin{exe}
  \ex \textbf{reconsiderado}
    \begin{xlist}
      \ex Ni siquiera me avisó que no iba. Jaja re piola, la mina,
          reconsiderada. Así estamos, hay gente para todo.
          (\url{https://www.facebook.com/karen.alma.90/posts/705446233192085})
      \ex Conicet: 47 aspirantes fueron reconsiderados e ingresarán a la carrera
          de investigador.
          (\url{https://www.clarin.com/sociedad/conicet-47-aspirantes-reconsiderados-ingresaran-carrera-investigador_0_r1NziPGxb.htm})
    \end{xlist}
  \ex \textbf{inutilizable}
    \begin{xlist}
      \ex Los celulares robados quedarán bloqueados y serán inutilizables.
          (\url{https://merloenlinea.com.ar/los-celulares-robados-quedaran-bloqueados-y-seran-inutilizables/})
      \ex El dirigible contratado por la empresa «La Serenísima» para fines
          publicitarios quedó inutilizable luego de un temporal en la localidad
          de Campo Andino, a 30 kilómetros de esta capital.
          (\url{https://www.lanacion.com.ar/sociedad/ahora-le-dicen-pinguino-porque-antes-volaba-nid170176})
    \end{xlist}
  \ex \textbf{desarmar}
    \begin{xlist}
      \ex Vasallo viene desarmando el discurso de la monogamia desde el 2013 y
          reflexiona sobre qué significa migrar de ese régimen.
          (\url{https://www.pagina12.com.ar/202982-los-malabares-del-deseo})
      \ex Al menos cinco ladrones rodearon y desarmaron este miércoles a
          Leonardo Sarmiento, un policía que trabajaba como encubierto en Villa
          Tranquila, Avellaneda\ldots{}
          (\url{https://www.infobae.com/sociedad/2018/11/02/menores-con-armas-de-juguete-amenazaron-a-una-periodista-en-vivo-al-policia-lo-queme-yo})
    \end{xlist}
\end{exe}

\begin{resuelto}
1. a. \forma{reconsiderado}: \sig{muy considerado o amable}.
\forma{considerar}\textsubscript{V} >
\forma{considerado}\textsubscript{Part/A}\footnote{El participio del verbo llega
a funcionar como adjetivo por un proceso derivativo conocido como conversión,
que estudiaremos en §\ref{sec:04-conversion}.} >
\forma{reconsiderado}\textsubscript{A};
[re-[[considera]\textsubscript{V}-d-o]\textsubscript{Part/A}]\textsubscript{A}; V >
Part/A > A. b. \sig{que se ha vuelto a considerar}:
\forma{considerar}\textsubscript{V} > \forma{reconsiderar}\textsubscript{V} >
\forma{reconsiderado}\textsubscript{V};
[[re-[considera]\textsubscript{V}]\textsubscript{V}-d-o]; V > V.
\textit{Aclaración}: en a., cuando el prefijo \forma{re-} se une a un A se
interpreta como intensivo (\sig{muy}); en b., cuando se une a un V, indica
repetición (\sig{volver a}).
\end{resuelto}
\end{ejercicio}

\section{Los procesos morfológicos}

Los morfemas de la lengua no se combinan unos con otros en forma aleatoria, sino
que lo hacen de manera sistemática y en un orden determinado. Cada uno de ellos
se incorpora a la palabra por medio de un proceso morfológico. Estos procesos se
dividen en dos grandes grupos: los \termino{de formación de palabras o
derivativos} y los \termino{flexionales}. Mientras que los primeros permiten
obtener palabras nuevas a partir de otras conocidas por el hablante, es decir,
crear nuevos \termino{lexemas}, por medio de los segundos se producen las
distintas \termino{variantes gramaticales} de un mismo lexema o unidad
vocabular de la lengua. Volviendo al ejemplo introducido en
§\ref{sec:04-segmentacion}, \forma{submarino}, \forma{submarina},
\forma{submarinos} y \forma{submarinas} son variantes gramaticales de un mismo
lexema adjetival, \forma{submarino}. Del mismo modo, \forma{mir-o, mir-á-s,
mir-a-ba-n, mir-a-r-ía-mos}, que cambian en los rasgos flexivos de tiempo,
persona y número, no son palabras distintas, sino formas que integran el
paradigma flexivo del mismo verbo \forma{mir-a-r}. Por contraste,
\forma{mir-a-dor}, que se forma añadiendo al verbo el sufijo derivativo
nominalizador \forma{-dor}, no es una variante del mismo lexema, sino otra
palabra diferente, que admite su propia variante flexiva plural
\forma{mir-a-dor-es}.

\subsection{La base}\label{sec:04-base}

La unidad sobre la cual se aplica un proceso morfológico se denomina
\termino{base}. Este concepto es relacional, puesto que la forma final de una
palabra puede obtenerse a partir de la aplicación ordenada de varios procesos
sucesivos. Retomando el ejemplo anterior, la base a partir de la cual se obtiene
\forma{mir-a-dor-es} es el tema del sustantivo \forma{mir-a-dor}, al cual se le
incorpora el sufijo flexivo \forma{-es}, que marca plural. Pero
\forma{mir-a-dor-} tiene como base, a su vez, el tema verbal \forma{mira-}, que
se convierte en nominal al adicionar el sufijo derivativo \forma{-dor}, que
forma un nuevo lexema.

\begin{exe}
  \ex\label{ex:04-mirador} \forma{mira(r)} > \forma{mira-dor} > \forma{mirador-es}
\end{exe}

\subsection{Productividad y transparencia}

Independientemente del hecho de que ambos grupos de procesos comparten el
recurso de la sufijación, poseen características bien diferenciadas, que
justifican una división del nivel morfológico en los dos ámbitos que
introdujimos en §\ref{sec:04-ramas}: la morfología léxica y la morfología
flexiva. Esta última se encuentra más próxima a la sintaxis y, por ello,
involucra procesos automáticos y obligatorios, es decir, que se pueden aplicar
sobre todas las palabras que poseen una determinada propiedad y no son
facultativos, sino que vienen determinados por los requisitos de buena formación
sintáctica de la oración en la que la palabra aparece. Así, todos los adjetivos
forman morfológicamente variantes de plural para concordar con los nombres a
los que modifican. La formación de palabras, en cambio, involucra procesos
optativos con distintos grados de productividad, que pueden no tener aplicación
sobre algunas bases. Por ejemplo, el sufijo \forma{-dad} produce sustantivos
abstractos como \forma{mal-dad, capac-i-dad, real-i-dad} a partir de adjetivos
como \forma{mal-o/a, capaz, real}, pero no se registra entre los hablantes una
forma como \juicio{*}{\forma{ternedad}} obtenida a partir de \forma{tiern-o/a}.
En cambio, se registra \forma{ternura}, obtenida de la incorporación del sufijo
\forma{-ura}, cercano semánticamente a \forma{-dad}, aunque menos productivo,
que también forma \forma{blancura}, \forma{dulzura} y \forma{verdura}. Podemos
decir, entonces, que la formación de un derivado abstracto en \forma{-ura}, en
el caso del adjetivo \forma{tierno}, bloquea o impide la formación de un
derivado abstracto en \forma{-dad}, que tendría el mismo significado.

Asimismo, por estar más cercana al léxico, esta área de la morfología tiende a
ser menos transparente semánticamente que la morfología flexiva y permite
interpretaciones más libres. Por caso, el significado de \forma{verdura}, en el
ejemplo anterior, no se interpreta habitualmente según su estructura
morfológica (\sig{cualidad de verde}), sino más bien como la designación de una
clase de alimentos, en muchos de los cuales predomina ese color. Ha ocurrido una
lexicalización por la que el significado de la palabra se ha vuelto más opaco.

\subsection{Un caso fronterizo: la morfología apreciativa}\label{sec:04-apreciativa}

La \termino{morfología apreciativa} constituye un caso especial, fronterizo entre
la derivación y la flexión. En español, comparte con esos dos procesos el
procedimiento de la sufijación, si bien, como señala \textcite{kornfeld2021},
también se expresa por medio de prefijos, como ocurre en usos ponderativos de
morfemas como \forma{re-} y otros análogos: \forma{re-leal, réquete-leal}. Al
igual que sucede con la morfología flexiva, la productividad de estos procesos
es muy alta (aunque no llega a ser automática). Por caso, casi todos los nombres
pueden formar diminutivos en \forma{-it-}, como \forma{cas-ita, airec-ito,
bosquec-ito}. Como se desprende de los ejemplos anteriores, los sufijos
apreciativos son selectivos respecto de la categoría de sus bases, pero no la
cambian, ni forman tampoco palabras nuevas al incorporarse. Presentan la
particularidad de que pueden repetirse, como en \forma{chiqu-it-it-o}. Sin
embargo, en consonancia con la formación de palabras, su aplicación no es
forzosa, no se encuentra determinada por la estructura sintáctica de la
oración, ya que estos sufijos no expresan información gramatical obligatoria. No
obstante, tampoco aportan significado léxico específico. Como se plantea en
Albano y Giammatteo (2000, p.~5),\pendiente{En la bibliografía figura Albano y
Giammatteo (2002)}

\begin{cita}
Por lo general, la semántica de los apreciativos se presenta relacionada con
parámetros tales como \sig{dimensión}, \sig{ponderación} y \sig{afectividad},
todos con la posibilidad de manifestar polaridad negativa tanto como positiva:
aumento/disminución (tamaño), valor/desvalor (ponderación), aprecio/desprecio
(afectividad).
\end{cita}

Es decir, los apreciativos se caracterizan por portar valores connotativos, que
codifican significados expresivos que quiere manifestar el hablante, lo cual
generalmente está muy vinculado al contexto en que se emiten tales expresiones.
Así, \textcite[128]{lang1997} reconoce lo siguiente:

\begin{cita}
\textit{-ito}, categorizado aisladamente como diminutivo y afectivo, puede ser
utilizado como peyorativo y despectivo en \textit{caudillo} >
\textit{caudillito}, o incluso como aumentativo en una construcción del tipo
\textit{comida calentita} (= comida muy caliente); \textit{-ote}, con claro
valor aumentativo, aparece en algunos contextos léxicos con un claro valor
diminutivo (cfr. \textit{isla} > \textit{islote}). La determinación del
significado concreto de los sufijos apreciativos en cada uno de los contextos
en que aparecen puede depender de factores externos a los sufijos mismos. El
valor peyorativo de \textit{caudillito} deriva del sentido negativo que se
aplica a la base \textit{caudillo}. O la intención del hablante, en función del
énfasis o el tono de voz, puede alterar el significado exacto ---\textit{¡qué
españolaza!} podría interpretarse como ponderación (\textit{¡qué españolaza más
guapa, más impresionante!}) o como desprecio (\textit{¡qué españolaza más
gorda, más fea!})---. Resultaría precipitado categorizar estos sufijos fuera de
su contexto.
\end{cita}

Distribucionalmente estos sufijos ocupan una posición intermedia entre la
morfología flexiva, que es más externa, y la léxica, que es más cercana a la
raíz, tal como se muestra en \forma{cas-er-it-o-s}, donde el sufijo apreciativo
\forma{-it-} se ubica entre el derivativo \forma{-er-}, que lo precede, y los
sufijos de flexión que indican género masculino \forma{-o} y número plural
\forma{-s}.\footnote{En algún sentido, esta generalización también parece
mantenerse para la prefijación, dado que, si bien en español no existen
prefijos flexivos, aquellos que participan del subsistema apreciativo resultan
más externos que los que aportan contenido léxico. Así, resulta aceptable
\textit{re-des-leal}, pero no *\textit{des-re-leal}.}

\begin{ejercicio}[El sufijo \textit{-ardo}]\label{ej:04-ardo}
Considere las siguientes imágenes, que incluyen palabras que han sido obtenidas
por un proceso que se popularizó en el léxico juvenil de las redes sociales en
los últimos años, y luego resuelva las consignas que se detallan más abajo:

\imagenpendiente{Imágenes del ejercicio: \textit{cafesardo, malardo, buenardo,
temardo, matesardo, quesardo}}

\textbf{Parte A}
\begin{consignas}
  \item Describa el proceso empleado en la formación de las palabras con el
        sufijo \forma{-ardo}, indicando la categoría de las bases y la del
        resultado.
  \item Realice una paráfrasis del significado del morfema en cuestión.
  \item Expanda el corpus con otros ejemplos de uso del mismo sufijo.
\end{consignas}

\begin{resuelto}[Parte A. Punto resuelto]
a) N\textsubscript{x}/A\textsubscript{y} + \forma{-ard-} >
N\textsubscript{x}/A\textsubscript{y}. (Los índices \textit{x} e \textit{y}
señalan que las bases no cambian su categoría).
\end{resuelto}

\textbf{Parte B}

Evalúe cuáles de las siguientes afirmaciones generalizan con mayor precisión las
conclusiones del análisis del ejercicio anterior.

\begin{consignas}
  \item Desde el punto de vista de la productividad del sufijo:
        \begin{exe}
          \ex El morfema se puede aplicar sobre cualquier base de una categoría,
              como los sufijos flexivos.
          \ex El morfema se puede aplicar sobre un conjunto acotado de bases,
              como los sufijos derivativos.
        \end{exe}
  \item Desde el punto de vista de la obligatoriedad:
        \begin{exe}
          \ex El proceso está determinado por la buena formación de la oración,
              como ocurre con la flexión.
          \ex El proceso es independiente de las condiciones de buena formación
              sintáctica, como ocurre con la formación de palabras.
        \end{exe}
  \item Desde el punto de vista del cambio categorial de la base:
        \begin{exe}
          \ex El proceso no cambia la categoría de la base.
          \ex El proceso cambia la categoría de la base.
        \end{exe}
  \item Desde el punto de vista del significado que aporta el sufijo a la
        formación:
        \begin{exe}
          \ex El morfema posee contenido léxico descriptivo, como los afijos
              derivativos.
          \ex El morfema posee contenido vinculado con las categorías
              lingüísticas, como los sufijos flexivos.
          \ex El morfema no tiene significado, como los interfijos.
          \ex El morfema tiene significado, pero su contenido no es descriptivo
              ni está vinculado con las categorías lingüísticas.
        \end{exe}
\end{consignas}

\textbf{Parte C}

Teniendo en cuenta los resultados obtenidos en su evaluación de los puntos
anteriores, ¿el sufijo \forma{-ardo} responde a las características de los
sufijos derivativos, a las de los flexivos o a las de los apreciativos?
\end{ejercicio}

\section{La formación de palabras}

Las lenguas distinguen entre palabras simples, que no provienen de otras, como
\forma{mar} o \forma{árbol}, y palabras complejas, formadas por procesos
morfológicos, como \forma{marítimo} y \forma{altamar}. La mayoría de las
palabras complejas se obtienen incorporando morfemas a una base, ya que la
morfología léxica es esencialmente aditiva o concatenativa. Según muestran los
ejemplos anteriores, los principales procesos de formación de palabras son la
\termino{derivación} y la \termino{composición}. La diferencia fundamental entre
ellos es la naturaleza de la unidad involucrada en el proceso morfológico:
mientras que en el caso de la derivación se añade un afijo, la composición une
dos o más palabras o temas. Mención aparte merece la \termino{parasíntesis}, en
la que ocurren de manera simultánea dos procesos, que pueden ser ambos de
naturaleza derivativa o uno derivativo y otro compositivo. A continuación,
presentaremos cada uno de estos procedimientos.

\subsection{La derivación}\label{sec:04-derivacion}

Una palabra derivada se obtiene por medio de la incorporación de un prefijo o un
sufijo a la base. Así, \forma{felic-i-dad} es una palabra derivada de
\forma{feliz} por medio del sufijo \forma{-dad} precedido del interfijo
\forma{-i-}; \forma{infeliz} también lo es, pero en este caso por acción del
prefijo \forma{in-}. La conversión y la formación regresiva se relacionan con
este proceso, con la diferencia de que no añaden segmentos morfológicos a la
base (v.~§\ref{sec:04-conversion} y §\ref{sec:04-regresiva}).

\begin{ejercicio}\label{ej:04-derivacion}
Considere el siguiente corpus y resuelva las consignas que se detallan a
continuación:

\begin{tabular}{@{}llll@{}}
  \forma{falsedades}  & \forma{famosas}        & \forma{chateaba}     & \forma{habilidad} \\
  \forma{bromear}     & \forma{deslealtad}     & \forma{cotidianidad} & \forma{amorosísimo} \\
  \forma{desinformó}  & \forma{bondades}       & \forma{desposteó}    & \forma{desagradables} \\
  \forma{grasosos}    & \forma{maldad}         & \forma{desamor}      & \forma{googlearon} \\
  \forma{ruidoso}     & \forma{disconformidad} & & \\
\end{tabular}

\textbf{Parte A}
\begin{consignas}
  \item Agrupe las palabras formadas por procesos comunes e indique de qué
        proceso se trata.
  \item Determine la categoría léxica de la base sobre la que se aplicó la regla
        y la categoría de la palabra resultante.
  \item Indique el valor semántico que los afijos aportan a la palabra.
  \item Realice la descripción fonológica del afijo incorporado a la estructura
        de la palabra.
\end{consignas}

\begin{resuelto}
Grupo 1. a) \forma{famosas, grasosos, amorosísimo, ruidoso}. Proceso común:
derivación con el sufijo \forma{-os-}. b) Categoría de la base: N. Categoría
resultante: A. c) Valor semántico: \sig{que tiene X}.\footnote{X = significado
de la base.} d) Descripción fonológica: \fonol{ˈos}.
\end{resuelto}

\textbf{Parte B}

A partir del análisis anterior, decida cuál de las siguientes generalizaciones
describe con mayor precisión lo observado en cuanto a la interacción entre los
morfemas derivativos y las categorías léxicas:

\begin{exe}
  \ex Todos los afijos derivativos son sensibles a la categoría de la base y
      cambian la clase o subclase de palabra.
  \ex Los sufijos derivativos son sensibles a la categoría de la base y
      típicamente cambian la clase o subclase de palabra; en cambio, los
      prefijos derivativos son típicamente insensibles a la categoría de la base
      y no cambian la clase o subclase de palabra.
  \ex Todos los sufijos derivativos son insensibles a la categoría de la base y
      no cambian la clase o subclase de palabra.
  \ex Los sufijos derivativos son insensibles a la categoría de la base y no
      cambian la clase o subclase de palabra; en cambio, los prefijos
      derivativos típicamente son sensibles a la categoría de la base y cambian
      la clase o subclase de palabra.
\end{exe}
\end{ejercicio}

\subsection{La composición}\label{sec:04-composicion}

Los compuestos se forman por la unión de dos o más palabras o temas. Se
denominan \termino{compuestos patrimoniales} aquellos que se forman con
palabras propias de la lengua, como en \forma{guardapolvo} (\forma{guarda} +
\forma{polvo}). Los \termino{compuestos grecolatinos}, en cambio, se producen
por la unión de los que se conocen como temas cultos
(v.~§\ref{sec:04-temascultos}), como en \forma{geógrafo} (\forma{geo} +
\forma{grafo}). También existen \termino{compuestos mixtos}, que se forman por
la unión de un tema culto y una palabra patrimonial, como \forma{hidroenergía} o
\forma{musicología}.

Según el grado de unión entre las bases, se reconocen \termino{compuestos
léxicos}, como \forma{sacacorchos, boquiabierto, agridulce}, y
\termino{compuestos sintagmáticos}, como \forma{ojo de buey, llave inglesa,
árabe-israelí}. En los compuestos léxicos las bases están plenamente soldadas y
unidas ortográficamente; en cambio, en los sintagmáticos, aunque tienen un
sentido unitario, los componentes no se escriben juntos e incluso pueden tener
una preposición entre ambos: \forma{máquina de coser, esmalte de uñas}.

La relación entre los elementos del compuesto puede ser de coordinación, como en
\forma{compraventa} y \forma{agridulce}, o de subordinación, como en
\forma{altiplanicie} o \forma{pelirrojo}. En los \termino{compuestos
coordinativos} ambos miembros pertenecen a la misma clase de palabra y pueden
parafrasearse por \forma{y} (\forma{agrio}\textsubscript{A} y
\forma{dulce}\textsubscript{A} a la vez); en cambio, en los \termino{compuestos
subordinativos} un miembro es modificador (\forma{pelirrojo}) o complemento
(\forma{marcapasos}) del otro.

\begin{ejercicio}\label{ej:04-composicion}
El siguiente corpus está integrado por distintos tipos de compuestos. Analícelos
y resuelva las consignas de abajo:

\begin{tabular}{@{}llll@{}}
  \forma{sacacorchos}      & \forma{sordomudo}      & \forma{aguardiente}      & \forma{lavarropas} \\
  \forma{socioeconómico}   & \forma{caradura}       & \forma{guardapolvos}     & \forma{blanquiceleste} \\
  \forma{pisapapeles}      & \forma{pelapapas}      & \forma{cazarrecompensas} & \forma{radioaficionado} \\
  \forma{dentoalveolar}    & \forma{matabichos}     & \forma{italoargentino}   & \forma{hierbabuena} \\
  \forma{cometierra}       & \forma{chupamedias}    & \forma{limpiaparabrisas} & \forma{agridulce} \\
\end{tabular}

\begin{consignas}
  \item Agrupe los que han sido obtenidos a partir de un esquema compositivo
        común.
  \item Segmente sus bases.
  \item A partir de uno de los ejemplos, realice la descripción estructural de
        cada esquema de palabra.
  \item Describa la regla que subyace a la formación de las palabras de cada
        grupo.
\end{consignas}

\begin{resuelto}
Grupo 1. a) \forma{aguardiente, caradura, radioaficionado, hierbabuena}. b)
\forma{agu(a)-ardiente; cara-dura; radio-aficionado; hierba-buena}. c)
[[agua]\textsubscript{N}[ardiente]\textsubscript{A}]\textsubscript{N}. d) N + A >
N.
\end{resuelto}
\end{ejercicio}

\begin{ejercicio}[Temas cultos]\label{ej:04-filos}
Considere el siguiente corpus integrado por elementos léxicos que comparten un
tema culto que proviene etimológicamente de la palabra del griego antiguo
\textit{philos} y resuelva las consignas de abajo:

\forma{filarmonía, filosofía, melófila, filonazi, hidrófilo, americanófila,
filología, tecnófilo, filántropo}

\begin{resuelto}[Ejemplo resuelto]
\forma{filarmonía}:
[[fil]\textsubscript{TC}[armonía]\textsubscript{N}]\textsubscript{N}; TC + N > N.
\end{resuelto}

\textbf{Parte A}

Realice la descripción estructural de cada palabra. Recuerde que los temas cultos
no pertenecen a ninguna clase de palabra, pero que las palabras que forman sí
están categorizadas.

\textbf{Parte B}

Compare el funcionamiento observado en los casos de arriba con el de los afijos
y el de las raíces y decida cuáles de las generalizaciones de abajo describen
con mayor precisión el comportamiento observado.

\begin{consignas}
  \item Con respecto a la distribución:
        \begin{exe}
          \ex al igual que los prefijos, los temas cultos son afijos que
              aparecen exclusivamente antes de la base.
          \ex al igual que los sufijos, los temas cultos son afijos que aparecen
              exclusivamente después de la base.
          \ex al igual que las raíces, los temas cultos son bases y pueden
              ocupar distintas posiciones en los esquemas compositivos.
          \ex los temas cultos son afijos, pero pueden ocupar distintas
              posiciones en los esquemas compositivos.
        \end{exe}
  \item Con respecto a la categoría gramatical:
        \begin{exe}
          \ex los temas cultos no poseen categoría gramatical propia ni
              determinan la categoría de la palabra resultante, al igual que los
              prefijos.
          \ex los temas cultos determinan la categoría de la palabra resultante
              al igual que los sufijos derivativos, pero no poseen categoría
              gramatical propia.
          \ex los temas cultos poseen categoría gramatical propia al igual que
              las raíces, pero no determinan la categoría de los compuestos.
          \ex los temas cultos poseen categoría gramatical propia y determinan
              la categoría de los compuestos al igual que las raíces.
        \end{exe}
\end{consignas}
\end{ejercicio}

\subsection{La parasíntesis}\label{sec:04-parasintesis}

Si bien se reconoce una tendencia a que cada paso de la historia derivacional de
las palabras esté asociado a un afijo derivativo, existen palabras que se forman
por medio de la incorporación simultánea de dos procesos, cada uno de los cuales
añade un morfema con contenido léxico a la base. Este fenómeno, que explica cómo
se obtiene la estructura de \forma{submarin-} sin haber pasado por
\forma{marin-} ni suponer un imposible \juicio{*}{\forma{submar}} (nuestro
interrogante en §\ref{sec:04-ramas}), se denomina \termino{parasíntesis}. La
peculiaridad de los procesos parasintéticos es que permiten suponer como paso
intermedio en la derivación una palabra inexistente.

En función del tipo de proceso que se aplica simultáneamente, existen dos tipos
de esquemas parasintéticos: \termino{por derivación} y \termino{por
composición}. El primero consiste en la combinación entre un proceso de
prefijación y otro de sufijación, como \forma{amontonar}, cuya base no es ni
\juicio{*}{\forma{montonar}} ni \juicio{*}{\forma{amontón}}, sino directamente el
sustantivo \forma{montón}, que incorpora en un mismo proceso el prefijo
\forma{a-} y el sufijo verbalizador formado por la vocal temática, al que se
añade la desinencia de infinitivo \forma{-r}, que es un morfema flexivo
\parencite[4702]{serranodolader1999}. Este primer tipo de parasíntesis en
español da lugar a verbos (\forma{enriquecer, agrandar, reciclar}) y a
adjetivos (\forma{anaranjado, aterrador, enrojecido}).

El segundo tipo de parasíntesis resulta de la formación de un derivado a partir
de un compuesto inexistente en la lengua al que se le añade un sufijo, como en
\forma{quinceañera}. Así, \forma{quince años}, que es una construcción
sintáctica y no una palabra compuesta existente en la lengua, igualmente actúa
como la base a la que se le aplica el sufijo \forma{-era} para formar el
derivado \forma{quinceañera}.

\begin{ejercicio}\label{ej:04-parasintesis}
Considere el siguiente corpus integrado por palabras parasintéticas y resuelva
las consignas de abajo:

\forma{desalmado, terraplanista, subterránea, mileuristas, salpimentar,
infumable}

\begin{consignas}
  \item Identifique en cada caso la base sobre la que se aplicó la parasíntesis.
  \item Realice la descripción estructural de la palabra y formalice los pasos
        de la derivación.
  \item Indique, en cada caso, de qué tipo de parasíntesis se trata (por
        derivación o por composición).
\end{consignas}

\begin{resuelto}
\forma{desalmado}. a) Base: \forma{alma-}. b)
[des-[alm]\textsubscript{N}-ad-o]\textsubscript{A}; \forma{des-} + N +
\forma{-ad-o} > A. c) Palabra parasintética por derivación.
\end{resuelto}
\end{ejercicio}

\section{Procesos menores de formación de palabras}

La morfología del español también dispone de mecanismos de creación léxica que
no agregan material fónico a la palabra. Entre ellos se pueden diferenciar, por
un lado, la \termino{conversión}, que no añade ningún afijo a la base, y la
\termino{formación regresiva} y los \termino{acortamientos}, que operan
sustrayendo formantes morfológicos a la base. Por el otro lado, otros procesos,
como los \termino{cruces léxicos} y los \termino{acrónimos}, directamente no
atienden a las fronteras morfológicas de la palabra y se explican a partir de
elisiones basadas en la estructura fonológica de la palabra. Por último, el
\termino{reanálisis} o \termino{etimología popular} es un proceso cuya
explicación recurre a la semántica, ya que lo que hace es reinterpretar el
significado de una parte de una palabra.

\subsection{La conversión}\label{sec:04-conversion}

Cuando una palabra cambia su categoría gramatical sin haber sufrido la
incorporación ni la sustracción de ningún segmento, se entiende que su
estructura morfológica sufrió una \termino{conversión}. La conversión se puede
considerar un proceso derivativo que, si bien no modifica los formantes
morfológicos, puede reanalizar, es decir, reinterpretar, alguno de los segmentos
de la palabra. En inglés es un proceso común que convierte sustantivos en
verbos: \forma{fish} (\sig{pez}) > \forma{(to) fish} (\sig{pescar});
\forma{house} (\sig{casa}) > \forma{(to) house} (\sig{alojar}). En español, por
ejemplo, los adverbios adjetivales, como \forma{rápido} en \forma{Llegué
rápido}, coinciden superficialmente con formas adjetivas flexionadas en
masculino singular. Sin embargo, su vocal final no puede considerarse un sufijo
flexivo porque, a diferencia de los adjetivos que flexionan en género y número,
los adverbios son palabras invariables; por tanto, la vocal final se ha
reanalizado como parte del tema adverbial. Así, en \exref{ex:04-rapido-a} y
\exref{ex:04-rapido-b}, \forma{rápido/rápida} son adjetivos predicativos que
concuerdan en género y número con sus respectivos sujetos (v. \capref{9}); en
cambio, \forma{rápido} en \exref{ex:04-rapidoadv} no manifiesta concordancia
porque es un adverbio equivalente a \forma{rápidamente}.

\begin{exe}
  \ex\label{ex:04-rapido}
    \begin{xlist}
      \ex\label{ex:04-rapido-a} Pedro es \forma{rápido} para hacer las cosas.
      \ex\label{ex:04-rapido-b} Margarita es \forma{rápida} para hacer las cosas.
    \end{xlist}
  \ex\label{ex:04-rapidoadv}
    \begin{xlist}
      \ex Él llegó \forma{rápido/rápidamente} a la estación.
      \ex Ella llegó \forma{rápido/rápidamente}.
    \end{xlist}
\end{exe}

Otro ejemplo de conversión en español puede considerarse el empleo de los
participios verbales \exref{ex:04-agotado-a} como adjetivos
\exref{ex:04-agotado-b} y \exref{ex:04-agotado-c}, y también el de los adjetivos
\exref{ex:04-terminal-a} como sustantivos \exref{ex:04-terminal-b}:

\begin{exe}
  \ex\label{ex:04-agotado}
    \begin{xlist}
      \ex\label{ex:04-agotado-a} Este trabajo nos \forma{ha agotado} a todos / a todas.
      \ex\label{ex:04-agotado-b} Juan quedó \forma{agotado} por el esfuerzo.
      \ex\label{ex:04-agotado-c} María quedó \forma{agotada} por el esfuerzo.
    \end{xlist}
  \ex\label{ex:04-terminal}
    \begin{xlist}
      \ex\label{ex:04-terminal-a} Llegamos temprano a la estación terminal.
      \ex\label{ex:04-terminal-b} La terminal estaba vacía.
    \end{xlist}
\end{exe}

Como vemos en \exref{ex:04-agotado-a}, el participio en los tiempos compuestos es
invariable, mientras que en el uso adjetivo debe concordar con el sustantivo al
que se refiere \exref{ex:04-agotado-b} y \exref{ex:04-agotado-c}. En
\exref{ex:04-terminal-a} tenemos un uso adjetivo de \forma{terminal} referido al
sustantivo \forma{estación}; en \exref{ex:04-terminal-b}, \forma{terminal} se
usa solo, precedido del determinante ---\forma{la terminal}---, por lo que
adquiere uso sustantivo.

\begin{ejercicio}\label{ej:04-conversion}
Las palabras que se destacan en los siguientes ejemplos, en ausencia de un
contexto sintáctico, podrían clasificarse en más de una
categoría.\footnote{Dado que el resultado de la conversión se vuelve perceptible
en el entorno sintáctico, este ejercicio entra en diálogo con los capítulos 5 y
11, sobre estructura del sintagma y perífrasis verbales. Recomendamos su
consulta para resolverlo.}

\begin{consignas}
  \item Segmente morfológicamente cada una de ellas, identificando raíz y tema.
  \item Analice los contextos sintácticos que corresponden a cada forma y
        permiten asignarla a una clase de palabra diferente en cada caso.
\end{consignas}

\begin{exe}
  \ex \textbf{cantar}
    \begin{xlist}
      \ex Aquí me pongo a cantar / Al compás de la vigüela, / Que el hombre que
          lo desvela / Una pena estrordinaria, / Como la ave solitaria / Con el
          cantar se consuela (José Hernández, \textit{El gaucho Martín Fierro},
          I).
      \ex Cantar de los cantares, el cual es de Salomón (\textit{Antiguo
          Testamento}, «Cantar de los cantares», 1:1).
    \end{xlist}
  \ex \textbf{comida}
    \begin{xlist}
      \ex El chef algoritmo: cómo es la comida hecha con inteligencia artificial
          que llegó a la Argentina.
          (\url{https://www.clarin.com/tecnologia/chef-algoritmo-comida-hecha-inteligencia-artificial-llego-argentina_0_znwOtd0bz.html})
      \ex Terror por una mujer que fue comida por un cocodrilo en Australia.
          (\url{https://www.minutouno.com/notas/3045423-terror-una-mujer-que-fue-comida-un-cocodrilo-australia/})
    \end{xlist}
  \ex \textbf{cursada}
    \begin{xlist}
      \ex Examen parcial es aquella evaluación que se toma en distintos momentos
          de la cursada de la materia, abarcando contenidos parciales.
          (\url{http://www.mdp.edu.ar/academica/ingreso/preguntas-uni.html})
      \ex [\ldots] a raíz de haber resultado infructuosa la notificación cursada
          en el domicilio que ahora denuncia la quejosa.
          (\url{https://www.abogados.com.ar/ratifican-validez-de-la-notificacion-del-traslado-de-la-demanda-practicada-en-el-domicilio-social-inscripto/16571})
    \end{xlist}
  \ex \textbf{privado}
    \begin{xlist}
      \ex Sin piso, el empleo privado cayó por sexto mes consecutivo en febrero.
          (\url{https://www.cronista.com/economiapolitica/Sin-piso-el-empleo-privado-cayo-por-sexto-mes-consecutivo-en-febrero-20190425-0040.html})
      \ex Artículo 61: El agente de la Administración Municipal no podrá ser
          privado de su empleo, ni objeto de sanciones disciplinarias sino por
          las causas y procedimientos determinados en este Estatuto.
          (\url{http://www.gob.gba.gov.ar/legislacion/legislacion/l-11757.html})
    \end{xlist}
  \ex \textbf{doctorando}
    \begin{xlist}
      \ex Me estoy doctorando, ergo, debo estar muy loca.
          (\url{http://lagataroma.blogspot.com/2017/02/me-estoy-doctorando-ergo-debo-estar-muy.html})
      \ex Un estudio asegura que los doctorandos son seis veces más propensos a
          desarrollar ansiedad o depresión en comparación con la población
          general.
          (\url{https://elpais.com/elpais/2018/03/15/ciencia/1521113964_993420.html})
    \end{xlist}
  \ex \textbf{cumbia}
    \begin{xlist}
      \ex \textoaincorporar{el fragmento de \textit{La luna y tú} (Ráfaga)}
      \ex Cumbia tu aldea y cumbiarás el mundo.
          (\url{https://www.pagina12.com.ar/200099-cumbia-tu-aldea-y-cumbiaras-el-mundo})
    \end{xlist}
\end{exe}

\begin{resuelto}
1. a) \forma{cant-a-r}. \forma{cantar}\textsubscript{a}: raíz \forma{cant-},
tema \forma{canta-}; \forma{cantar}\textsubscript{b}: \forma{cant-ar(-es)}, raíz
\forma{cant-}, tema \forma{cantar-}. b) \forma{cantar}\textsubscript{a}: se
trata de un verbo porque integra una perífrasis verbal (v. \capref{11}):
\forma{ponerse a cantar} (= \sig{empezar a cantar});
\forma{cantar}\textsubscript{b}: se trata de un sustantivo porque va precedido
del determinante \forma{el} y puede pluralizarse: \forma{los cantar-es}.
\end{resuelto}
\end{ejercicio}

\subsection{La formación regresiva}\label{sec:04-regresiva}

Hay algunas palabras que, contra lo que podríamos suponer, no se han formado por
un incremento afijal, sino por la sustracción de un morfema de la base, en un
proceso conocido como \termino{formación regresiva}. Así, por ejemplo, en la
historia de la lengua está atestiguado que, contra lo esperado, no es
\forma{deslizar} la forma que proviene de \forma{desliz}, sino a la inversa: el
nombre \forma{desliz} se formó a partir del verbo \forma{deslizar}. Lo mismo ha
sucedido con el verbo \forma{legislar}, que deriva de \forma{legislador} y no al
revés, es decir que para su formación se eliminó el sufijo \forma{-dor} de la
palabra base.

\begin{ejercicio}\label{ej:04-regresiva}
Considere los siguientes enunciados, que contienen palabras neológicas acuñadas
para un ámbito específico o para causar un efecto publicitario o humorístico, y
resuelva las siguientes consignas:

\begin{consignas}
  \item Identifique los temas que funcionan como base y resultado del proceso
        con sus categorías y el segmento morfológico afectado.
  \item Formule como una regla el proceso morfológico involucrado en la
        formación del resultado.
  \item Realice una paráfrasis del significado de la palabra resultante del
        proceso.
\end{consignas}

\begin{exe}
  \ex \textbf{televisar}: El canal que empezará a televisar la NBA en la
      Argentina y cómo se verá la jornada inaugural.
      (\url{https://www.lavoz.com.ar/deportes/basquet/el-canal-que-empezara-a-televisar-la-nba-en-la-argentina-y-como-se-vera-la-jornada-inaugural})
  \ex \textbf{espectar}: El modo espectador es casi igual al modo creativo,
      solo que, como el nombre lo dice, te da la posibilidad de espectar un mundo
      pero no interactuar con este\ldots{}
      (\url{https://aminoapps.com/c/mcft-s/page/item/espectador/KW0u_KIRYJlRb8pNnmb26bqbKj3QRBz})
  \ex \textbf{alfabeto}: La diferencia entre un alfabeto funcional y un
      analfabeto funcional radica en que el alfabeto funcional pondrá en
      práctica competencias, conocimientos y habilidades adquiridas durante su
      vida académica\ldots{}
      (\url{http://transformacion-educativa.com/2do-congreso/ponencias/Eje-1/L2-102.html})
  \ex \textbf{antitranspirar}\pendiente{Falta el ejemplo en el original}
\end{exe}

\begin{resuelto}
1. a) Base: \forma{televisión}. Categoría: N. Resultado: \forma{televis(a)-}.
Categoría: V. Segmento afectado: \forma{-sión}. b) \forma{televisar} <
\forma{televisión}: N -- \forma{-sión} > V. c) Significado: \sig{transmitir
imágenes a distancia}.
\end{resuelto}
\end{ejercicio}

\subsection{El acortamiento}

Una palabra acortada resulta de la sustracción de material fonológico de la
base. Como resultado de este proceso, se conserva o bien su parte inicial,
proceso que se conoce como \termino{apócope}, o bien la parte final, lo que se
denomina \termino{aféresis}. Son apócopes de uso cotidiano: \forma{celular} >
\forma{celu}; \forma{señorita} (= \sig{maestra}) > \forma{seño}/\forma{se};
\forma{bicicleta} > \forma{bici}. Como ejemplos de aféresis podemos mencionar:
\forma{violonchelo} > \forma{chelo} y \forma{autobús} > \forma{bus}. Un ejemplo
característico de este proceso lo constituyen los hipocorísticos, es decir, los
acortamientos de los nombres propios: \forma{Isa} por \forma{Isabel};
\forma{Maxi} por \forma{Maximiliano}. Muchos nombres permiten tanto el
acortamiento por apócope como por aféresis: \forma{Fer} y \forma{Nando}, por
\forma{Fernando}; \forma{Guada} y \forma{Lupe}, por \forma{Guadalupe}.

\begin{ejercicio}\label{ej:04-acortamiento}
Observe el siguiente corpus integrado por palabras acortadas y a partir de ellas
resuelva lo pedido más abajo:

\forma{cheto, profe, facu, tano, ñora, auto, moto, compu, ñero/ñeri, tamos,
taxi}

\begin{consignas}
  \item Agrupe las que se encuentran formadas por apócope o aféresis.
  \item Identifique sus bases y determine su categoría y la del resultado.
  \item Señale qué principio fónico se observa en la selección de los segmentos
        conservados.
  \item Evalúe qué modificación semántica produce la aplicación de estos
        procesos.
\end{consignas}

\begin{resuelto}
Grupo 1. a) Acortamientos por aféresis: \forma{cheto, tano, ñora, ñero/ñeri,
tamos}. b) \forma{cheto}\textsubscript{A} < \forma{concheto}\textsubscript{A};
\forma{tano}\textsubscript{A} < \forma{napolitano}\textsubscript{A};
\forma{ñora}\textsubscript{N} < \forma{señora}\textsubscript{N};
\forma{ñero/ñeri}\textsubscript{N} < \forma{compañero}\textsubscript{N};
\forma{tamos}\textsubscript{V} < \forma{estamos}\textsubscript{V}. c) Se
conservan las últimas dos sílabas de la base. La palabra resultante tiene
acentuación grave. d) El acortamiento por aféresis no produce cambio en el
significado de la base.
\end{resuelto}
\end{ejercicio}

\subsection{El cruce léxico o \textit{blending}}

Los \textit{blends} o «palabras macedonia» están formados a partir de los
extremos fónicos de otras palabras que se combinan entre sí. Estos extremos no
coinciden con morfemas, sino que pueden ser sílabas u otras secuencias de
segmentos tomadas del comienzo o del final de la palabra. Frecuentemente, el
punto de encuentro entre las bases coincide fónicamente. Por ejemplo, la palabra
\forma{estanflación} se obtiene a partir de fragmentos de las bases
\forma{estancamiento} e \forma{inflación} fusionando, respectivamente, los
extremos anterior (\forma{estan-}) y posterior (\forma{-flación}) al primer
sonido nasal que comparten (\forma{-n-}):
[\forma{estan(camiento)}]\textsubscript{N} +
[\forma{(i)nflación}]\textsubscript{N} >
[\forma{estanflación}]\textsubscript{N}. Por la naturaleza de las unidades
involucradas, el cruce léxico es un proceso morfológico que forma palabras
compuestas, cuya unión involucra reajustes fonológicos.

\begin{ejercicio}\label{ej:04-cruces}
El siguiente corpus está integrado por cruces léxicos de distinto tipo. Luego de
observar sus diferencias, realice las siguientes consignas:

\begin{consignas}
  \item Identifique las bases a partir de las cuales se formó cada compuesto.
  \item Determine la categoría de ambas bases y la del resultado.
  \item Analice cómo han sido seleccionados los segmentos de las bases
        utilizados en el compuesto.
\end{consignas}

\begin{exe}
  \ex Un \textbf{módem} es el dispositivo que permite a tus computadoras,
      teléfonos inteligentes, tabletas y otros dispositivos conectarse a
      Internet. (\url{https://www.lenovo.com/ar/es/glosario/que-es-modem/})
  \ex \textbf{Conmebol}
  \ex \textbf{Mercosur}
  \ex \textbf{vitalitud}
  \ex \textbf{deshamburguesar}
  \ex \textbf{artivismo}
  \ex \textbf{juernes}
\end{exe}

\imagenpendiente{Imágenes de los ítems 2 a 7}

\begin{resuelto}
1. \forma{módem}. a) Bases: \forma{modulador} + \forma{demodulador}. b)
Categoría de las bases: N + N. Categoría del resultado: N. c) Se seleccionaron
los primeros tres segmentos sonoros de cada una de las palabras, aprovechando el
fonema consonántico común: [mod(ulador) + dem(odulador)].
\end{resuelto}
\end{ejercicio}

\subsection{La acronimia}

La \termino{acronimia} resulta de la lexicalización, o conversión en palabra, de
una sigla. Las \termino{siglas} son formaciones que se producen a partir de las
iniciales de las palabras de clase mayor (sustantivos, adjetivos, verbos) de una
denominación compuesta por varias palabras. En algunos casos, también pueden
estar integradas por más de uno de los sonidos iniciales de cada palabra
(\forma{RENAPER}, por \forma{Registro Nacional de las Personas}) o incluir
palabras de clases menores (\forma{ADEBA}, por \forma{Asociación de Bancos de
Buenos Aires}). También existen siglas que resultan de una combinación entre
ambas estrategias, como \forma{EDENOR} (\forma{Empresa Distribuidora Norte}).
Pueden pronunciarse fonémicamente, es decir, de acuerdo con la realización
habitual de cada uno de sus segmentos sonoros (por ejemplo, \forma{ARCA}, para
\forma{Agencia de Recaudación y Control Aduanero}), o por medio de la mención de
los nombres de los grafemas que las integran (\forma{GNC}, pronunciado
\fonol{xeneˈse}, para \forma{Gas Natural Comprimido}).\pendiente{El original dice
«Gas Natural Licuado»}

Todas las siglas mencionadas anteriormente se encuentran lexicalizadas, por lo
que se emplean como si fueran palabras de la lengua; así, la \forma{Universidad
de Buenos Aires} es popularmente conocida como la \forma{UBA} y la
\forma{Organización de las Naciones Unidas}, como la \forma{ONU}. Las siglas
lexicalizadas, que pasan a funcionar en la lengua como palabras, incluso
permiten ser derivadas: por ejemplo, de \forma{SIDA} (\forma{síndrome de
inmunodeficiencia adquirida}) se ha formado \forma{sidoso} (\sig{la persona que
padece este síndrome}) y de \forma{OVNI} (\forma{objeto volador no
identificado}) se ha derivado \forma{ovnilogía} (disciplina que estudia los
ovnis).\footnote{También se emplea \textit{ufología}, derivado de la sigla en
inglés UFO (\textit{unidentified flying object}).}

\begin{ejercicio}\label{ej:04-acronimos}
El siguiente corpus está integrado por distintos acrónimos. En cada caso:

\begin{consignas}
  \item Identifique sus bases y determine su categoría.
  \item Determine la categoría de la construcción resultante.
  \item Observe si existen marcas morfológicas que evidencien la lexicalización
        del acrónimo.
  \item Evalúe qué modificación semántica produce la aplicación del proceso.
\end{consignas}

\begin{exe}
  \ex \textbf{TICs}: Una mirada institucional sobre la incorporación de las TICs.
      (\url{https://twitter.com/andynsane/status/1085360342513958914?lang=es})
  \ex \textbf{RAE}: La RAE vela, desde 1713, por el buen uso y la unidad de la
      lengua española en el mundo. (\url{https://www.instagram.com/laraeinforma/})
  \ex \textbf{ovnis}: La cantidad de personas que cree que los avistamientos de
      ovnis ofrecen una prueba probable de vida extraterrestre aumentó del 20\,\%
      en 1996 al 34\,\% en 2022.
      (\url{https://www.lanacion.com.ar/lifestyle/en-las-redes/la-peligrosa-idea-de-que-extraterrestres-visitaron-la-tierra-esta-cada-vez-mas-extendida-nid09092024/})
  \ex \textbf{FIFA / AFA}: La FIFA advirtió «severas sanciones» si el Gobierno
      interviene la AFA.\pendiente{El enlace del original repite el del ítem 1}
  \ex \textbf{UNA}\pendiente{Falta el ejemplo en el original}
\end{exe}

\begin{resuelto}
1. a) \forma{TICs}: tecnología\textsubscript{N}, información\textsubscript{N},
comunicación\textsubscript{N}. b) Categoría: N. c) El segmento final corresponde
con la pluralización, esperable para un nombre, si bien no se manifiesta por
medio del alomorfo correspondiente a la terminación del tema (\forma{-es}). d)
El acrónimo equivale referencialmente a la frase que abrevia.
\end{resuelto}
\end{ejercicio}

\subsection{Reanálisis o etimología popular}

El \termino{reanálisis} morfémico o simplemente reanálisis
\parencite{escandell2011}, también conocido como \termino{etimología popular}, es
el proceso por el que el hablante analiza las partes de una palabra y las
reinterpreta asignándoles significados que no tenían originalmente. Así, en
inglés, la palabra \forma{hamburger} originalmente era el gentilicio
correspondiente a los habitantes de la ciudad alemana de Hamburgo, pero la
primera parte se confundió con \forma{ham} (\sig{jamón}) y la palabra se
reinterpretó como un compuesto [[ham] [burger]], lo que dio origen a otros
compuestos, como \forma{cheeseburger} (\sig{hamburguesa de queso}),
\forma{fishburger} (\sig{hamburguesa de pescado}), \forma{beefburger}
(\sig{hamburguesa de carne}) y también a la creación, por formación regresiva,
de la palabra \forma{burger}, como denominación general de cualquier tipo de
hamburguesa. Otro caso muy conocido de reanálisis es el de la palabra
\forma{gate} (lit. \sig{portón}) en \forma{Watergate} (lit. \sig{portón de
agua}), nombre de las oficinas en las que, a principios de los 70, se produjo un
escándalo de espionaje que obligó a dimitir al presidente de los EE.\,UU.,
Richard Nixon. El segundo elemento del compuesto, \forma{-gate}, comenzó a
utilizarse en otras palabras con el nuevo significado de \sig{asunto
escandaloso}: \forma{Irangate} (venta de armas a Irán), \forma{sexgate}
(\textit{affaire} sexual del presidente Clinton en la Casa Blanca). El mecanismo
también ha sido adoptado en español: \forma{Yomagate} (caso de narcotráfico y
corrupción que involucraba a la familia Yoma), \forma{FIFAgate} (escándalo de
corrupción en la FIFA), etc.

El reanálisis también se puede producir con fines humorísticos, como se muestra
en la viñeta de la figura~\ref{fig:04-afortunado}, en la que, en el adjetivo
parasintético \forma{a-fortun-ad-o} (\sig{que tiene fortuna o suerte}), el
prefijo inicial se reinterpreta como si se tratara del prefijo negativo
\forma{a-} (\forma{a-moral, a-normal}, etc.) aplicado a \forma{fortunado},
produciendo el significado opuesto de \sig{desprovisto de fortuna}.

\figurapendiente{Reanálisis humorístico de \textit{afortunado}}{fig:04-afortunado}

\begin{ejercicio}[\textit{Siempre x, nunca in-x}]\label{ej:04-siempre}
La construcción sintáctica \forma{siempre x, nunca in-x}, de uso frecuente en las
redes sociales, se emplea para oponer una palabra al que se considera su
antónimo, para reafirmarla como propiedad caracterizadora del que la enuncia.
Esta construcción se popularizó como meme en las redes sociales a raíz de un tuit
de la modelo chilena Nicole Moreno, que tenía como una de sus frases
características \forma{Siempre fuerte, nunca infuerte} (\textit{sic}: «siempre
foerte nunca infoerte»). Por la productividad del proceso morfológico
subyacente, el esquema se puso de moda y produjo una catarata de neologismos
lúdicos y efímeros, como el de la propia Moreno. Más interesante aún, con la
misma calidad de fugaces y lúdicas, se acuñaron palabras como las que se
destacan en los ejemplos que, junto con otros de origen diverso, se reproducen a
continuación. A partir de ellas, resuelva las siguientes consignas:

\begin{consignas}
  \item Determine la categoría de las bases sobre las que se aplica esta
        operación y la de las palabras resultantes.
  \item Identifique el segmento que es reinterpretado en la formación de estas
        palabras e indique el significado que se le asigna.
  \item Explique qué particularidad presenta \forma{imberbe} > \forma{berbe}.
\end{consignas}

\begin{exe}
  \ex \textbf{bécil}: Siempre bécil, nunca imbécil.
      (\url{https://x.com/PabloGa27864933})
  \ex \textbf{pecable}: Siempre pecable, nunca impecable.
      (\url{https://x.com/teseodeseo_?lang=ar})
  \ex \textbf{berbe}: Siempre berbe, nunca imberbe.
      (\url{https://twitter.com/filozoofico})
  \ex \textbf{teligente}\pendiente{Falta el ejemplo en el original}
\end{exe}

\begin{resuelto}
a) 1. \forma{bécil}. Base: \forma{imbécil}. Categoría: A. Resultado:
\forma{bécil-}. Categoría: A. b) \forma{bécil} designa la propiedad opuesta a
\forma{imbécil}.\pendiente{El original agrega «c) Ejemplo: alfabeto: A --- an- >
A», que corresponde al ejercicio~\ref{ej:04-regresiva}}
\end{resuelto}
\end{ejercicio}

\begin{ejercicio}\label{ej:04-reanalisis}
¿Cómo se podría explicar la formación de las siguientes palabras? ¿Se trata de
casos de reanálisis? Justifique su respuesta teniendo en cuenta la palabra de la
que derivan y la interpretación del significado.

\begin{exe}
  \ex \textbf{flamable}\imagenpendiente[14mm]{Imagen del ítem 1}
  \ex \textsc{El periodista}.--- Mire, señor militar: me dicen que antes era
      dictadura y ahora \textbf{dictablanda}. ¿Qué les debo contestar?
      (\textit{Diario El Sol de España}, 12/03/1930;
      \url{https://cvc.cervantes.es/el_rinconete/anteriores/septiembre_12/17092012_02.htm})
  \ex {[\ldots]} \textit{Episodio IV: Una nueva esperanza} fue la primera
      película de Star Wars, el inicio de la saga como la conocemos. Mientras
      que los episodios V y VI se convirtieron en secuelas, los episodios I, II
      y III fueron las \textbf{precuelas}.
\end{exe}

\begin{resuelto}
1. \forma{flamable}. Es un reanálisis o reinterpretación humorística del prefijo
locativo \forma{in-}, presente en \forma{in-tubar} o \forma{in-filtrar}, como si
se tratara del prefijo negativo homónimo (\forma{in-sano, in-moral}). De ahí
que, paradójicamente, el significado de la palabra resultante no cambie:
\sig{que puede producir llamas} (quemarse).
\end{resuelto}
\end{ejercicio}

\section{Las variantes morfológicas de las palabras}\label{sec:04-variantes}

Uno de los modos posibles de clasificación de las palabras que forman el léxico
de la lengua permite separarlas en categorías variables e invariables. Esto
quiere decir que hay clases de palabras que oponen de forma morfológica
manifiesta determinados contenidos semánticos altamente gramaticalizados, como
la cantidad ---\forma{perro/perros}--- o el tiempo de ocurrencia de un suceso
---\forma{canté/canto/cantaré}---. Las manifestaciones lingüísticas de estos
significados generales se conocen como categorías lingüísticas. Las que se
reconocen para el español son las siguientes: género, número, persona, caso,
tiempo, modo, aspecto y voz (aunque esta última no tiene manifestación
morfológica en español). La expresión de estas categorías, hecha la excepción de
la voz, se da a través de la flexión, que Fábregas (2015, p.~175)\pendiente{En la
bibliografía figura Fábregas (2013)} define como el «conjunto de operaciones que
nos proporciona la forma que una palabra debe adoptar dado cierto contexto
sintáctico». Efectivamente, la flexión es relevante para la sintaxis porque en
algunos casos expresa de manera abierta la existencia de determinadas relaciones
gramaticales ---las distintas formas de la concordancia--- y es un factor
esencial en la buena formación de la oración. Así, ni \juicio{*}{\forma{el gato
blanca}} ni \juicio{*}{\forma{los gato blancas}} son sintagmas bien formados, ya
que, según se verá en el próximo capítulo, tanto el artículo como el adjetivo
deben concordar en género y número con el nombre núcleo del sintagma, de modo
que la forma gramatical sería \forma{el gato blanco}. La morfología flexiva del
español es fundamentalmente aditiva: se realiza mediante la anexión de sufijos a
la base de la palabra.

\subsection{Las categorías gramaticales y su manifestación en la morfología}\label{sec:04-categorias}

Las categorías lingüísticas son moldes a través de los cuales la mente
conceptualiza el mundo \parencite[6]{lakoff1987}. Se expresan, como hemos
visto, morfológicamente, pero también pueden hacerlo léxicamente. Por ejemplo,
para el género, además de contrastes como \forma{perr-o/perr-a}, que se
manifiestan en la oposición entre las desinencias \forma{-o} y \forma{-a},
también hay oposiciones léxicas (mediante cambio de palabra) ---\forma{vaca/toro}---
o sintácticas (mediante el cambio del determinante) ---\forma{el dentista / la
dentista}---.

Las categorías gramaticalizan ---es decir, manifiestan mediante los recursos
gramaticales de la lengua--- aspectos que los hablantes extraen de la situación
extralingüística. En \textcite{giammatteo2018} se propone una clasificación que
parte de esta vinculación de las categorías con la situación comunicativa. La
clasificación da cuenta de una gradación que va de aquellas categorías
fuertemente enraizadas en la situación comunicativa hasta las más abstractas y
aparentemente menos vinculadas al contexto de situación. En función de su grado
de relación con el entorno en el que se pone en uso la lengua, las categorías se
organizan en tres grupos: las de la situación comunicativa, las del evento y las
de la escena oracional.

En primer lugar, las \termino{categorías de la situación comunicativa} son las
que están estrechamente relacionadas con los participantes del coloquio y el
marco en el que se produce el intercambio lingüístico. En este grupo se incluyen
las categorías de persona y tiempo, que se consideran categorías deícticas
porque establecen «un tipo de vínculo referencial entre ciertas unidades o
expresiones lingüísticas y aquello que representan en el mundo o en el universo
del discurso» \parencite[932]{eguren1999}. La tercera categoría de este grupo,
el modo, si bien no es deíctica, establece una relación indirecta con la
situación comunicativa, ya que tiene que ver con el hablante y con la actitud
que asume respecto de lo que dice.

\begin{itemize}
  \item La \termino{persona} señala a los participantes de la comunicación: el
        emisor (\forma{yo} (sg.) / \forma{nosotros/as} (pl.)); el destinatario
        (\forma{vos / tú / Ud.} (sg.) // \forma{Uds. / vosotros/as}
        (pl.))\footnote{La forma \textit{vosotros/as} solo se usa en el español
        peninsular, mientras que \textit{tú} se utiliza en España y en otras
        regiones de Latinoamérica. \textit{Vos} es la forma característica del
        Río de la Plata, pero se emplea también en algunas otras partes de
        Latinoamérica. Mientras \textit{tú} y \textit{vos} indican confianza,
        \textit{Ud.} se emplea en el trato formal.} y un tercer participante que
        no forma parte de la relación de alocución (\forma{él / ella} (sg.) //
        \forma{ellos / ellas} (pl.)). Este tercer participante puede o no estar
        en la situación comunicativa y puede ser cualquier tipo de entidad,
        \sig{animada} o no. Por eso a la tercera persona también se la conoce
        como la «no persona» y es la que se emplea en todos los usos vinculados
        con la impersonalidad: \forma{Mañana llueve / Se habla mucho y se hace
        poco}. Además de en los pronombres, la persona se manifiesta en la
        desinencia del verbo para concordar con la persona del sujeto:
        \forma{Yo cant-o / Él cant-a}.
  \item El \termino{tiempo} relaciona un evento con un punto o intervalo
        temporal respecto del cual el hecho se sitúa. Este momento temporal
        puede ser extralingüístico y ubicarse directamente como anterior,
        simultáneo o posterior respecto del momento de enunciación, como en
        \forma{Ayer fui al cine}, que indica un evento sucedido con anterioridad
        al momento de la emisión; o bien puede orientarse indirectamente a partir
        de la referencia proporcionada por otro evento, como en \forma{Llegó}
        (Ev\textsubscript{principal}) \forma{cuando todos se habían ido}
        (Ev\textsubscript{2 anterior}).
  \item El \termino{modo} codifica la relación entre el enunciador y lo que
        dice, tanto en términos de la actitud que adopta en la enunciación
        (aseverar, preguntar, ordenar, etc.) como en cuanto al grado de realidad
        que le adjudica al evento (probable, obligatorio, etc.).
\end{itemize}

Tiempo y modo son dos categorías que se manifiestan en la desinencia verbal,
muchas veces en forma de amalgama (v.~§\ref{sec:04-vt}). Así, por ejemplo, en
\forma{cant-a-ba-mos} reconocemos la raíz \forma{cant-}, la vocal temática
\forma{-a-}, la amalgama de tiempo, modo y aspecto que se conoce como TAM y la
amalgama de persona y número \forma{-mos}.

En segundo lugar, las \termino{categorías del evento} son tres: el género, el
número y el aspecto. Estas categorías nos permiten clasificar, cuantificar o
delimitar las entidades de las que hablamos y los sucesos en los que intervienen
y que constituyen el tema de nuestras conversaciones.

\begin{itemize}
  \item El \termino{género} funciona esencialmente como un clasificador de las
        entidades, puesto que divide a los nombres en clases (femenino,
        masculino, neutro u otras, según las lenguas). En español solo se
        interpreta de forma sistemática, en términos de oposición de sexo, para
        los nombres animados (de persona y de animales superiores o domésticos):
        \forma{abogado/abogada; león/leona; gallo/gallina}; en cambio, en los
        nombres de entidades inanimadas resulta inmotivado o inherente (propio
        de la palabra): \forma{el árbol, la piedra}, etc. Aunque no todos los
        sustantivos terminan en vocal, existe una tendencia a que el morfema
        \fonol{-a} funcione como marca de femenino y \fonol{-o} marque el
        masculino, aunque hay excepciones como \forma{(la) mano} o \forma{(el)
        clima}. Además, la vocal \fonol{e} se distribuye complementariamente con
        el morfema \fonol{-o} como marca de masculino (\forma{elefante, jefe}),
        aunque también funciona como marca de ambos géneros en los adjetivos
        (\forma{elegante, amable}). La marcación de género femenino puede
        resultar suplida por otras desinencias (\forma{gallo/gall-ina,
        conde/cond-esa, actor/actr-iz}).
  \item El \termino{número} refiere a la cuantificación de las entidades (una o
        más de una) en los nombres ---\forma{gato/gatos}--- y a la de las
        personas intervinientes en el suceso, en el verbo ---(\forma{él})
        \forma{corre} / (\forma{ellos}) \forma{corren}---. En el verbo se
        manifiesta para concordar con el nombre núcleo del sujeto, junto con la
        categoría de persona: \forma{Yo cant-o / Nosotros cant-a-mos}.
  \item El \termino{aspecto} sirve tanto para expresar la forma interna del
        desarrollo de los eventos (si son dinámicos o estáticos, si poseen o no
        duración, si cuentan con una delimitación inherente o si pueden
        repetirse), como también para hacer referencia a sus alternativas de
        focalización (el punto de vista desde el que se los presenta:
        completados, en desarrollo o bien en alguna de sus fases, como el inicio
        o la conclusión). La principal, y tal vez la más extendida, oposición
        aspectual en las lenguas es la que distingue entre el aspecto
        imperfectivo, que muestra al evento en desarrollo, y el aspecto
        perfectivo, que lo concibe desde afuera, como un todo limitado. Este
        contraste se puede ilustrar para el español con la oposición entre
        \forma{cantó}, que indica un hecho puntual, y \forma{cantaba}, que
        señala un evento durativo.\footnote{Las nociones relativas al aspecto se
        amplían en el \capref{8}, donde se trata la estructura eventiva. También
        se puede consultar \textcite{debin2018}.}
\end{itemize}

Finalmente, las \termino{categorías de la escena oracional}, el caso y la voz o
diátesis (que no aparece en la flexión en español), se encargan de reproducir
dentro de la oración las relaciones que se establecen entre el evento y los
participantes que intervienen en él, manifestados, respectivamente, por el verbo
y sus argumentos.\footnote{Estas cuestiones se desarrollan en el \capref{7},
dedicado a la estructura argumental del verbo y a las funciones sintácticas.}

\begin{itemize}
  \item El \termino{caso} indica el tipo de relación que vincula a los sintagmas
        nominales, que representan los argumentos del verbo, con el predicado.
        Aunque en las lenguas clásicas, como el griego y el latín, y también en
        lenguas modernas, como el ruso y el alemán, estas diferencias se
        manifiestan morfológicamente, en español no tenemos una marca
        morfológica para los diferentes casos que indican las funciones
        sintácticas de los argumentos del verbo, excepto cuando se trata de
        algún pronombre. Los pronombres personales manifiestan caso nominativo,
        que es el que corresponde a la función sujeto ---\forma{yo}---, caso
        objetivo, para las funciones de COD y COI ---\forma{me}--- y caso oblicuo
        o terminal, para cuando siguen a una preposición ---(\forma{para})
        \forma{mí}---.\footnote{Este tema se amplía en el \capref{10}, §2.}
  \item La \termino{diátesis} es el nombre que reciben cada una de las
        construcciones gramaticales que permiten expresar de manera diversa los
        argumentos de un verbo; por ejemplo, la alternancia que se da entre la
        voz activa ---\forma{Yo leo un libro}--- y la pasiva ---\forma{Un libro
        es leído por mí}---. Como se muestra en el ejemplo, la diátesis no tiene
        manifestación morfológica en español, sino que se expresa mediante una
        construcción sintáctica.\footnote{La perífrasis de voz pasiva se trata
        en el \capref{11} y la pasiva cuasirrefleja, en el \capref{10}.}
\end{itemize}

\begin{ejercicio}\label{ej:04-genero}
Existe una percepción extendida entre los hablantes del español según la cual en
el entorno nominal la vocal \fonol{a} en posición final de palabra significa
\sig{género femenino} y la vocal \fonol{o} en esa misma posición se interpreta
como \sig{género masculino}. A partir del siguiente corpus de sintagmas
nominales del español resuelva lo pedido más abajo:

\begin{tabular}{@{}lllll@{}}
  1. la secretaria & 2. el sartén   & 3. la gata    & 4. la cebra    & 5. el naranjo \\
  6. la cuchilla   & 7. la araña    & 8. la mano    & 9. el clima    & 10. el alcalde \\
  11. el amante    & 12. el tilde   & 13. la bolsa  & 14. el calor   & \\
\end{tabular}

\textbf{Parte A}

Cambie el género gramatical de los nombres.

\begin{resuelto}
\forma{la secretari-a / el secretari-o}.
\end{resuelto}

\textbf{Parte B}

Decida cuáles de las siguientes generalizaciones se adecuan mejor a lo
observado:

\begin{consignas}
  \item Con respecto a la disponibilidad de la oposición en género gramatical:
        \begin{exe}
          \ex Todos los sustantivos expresan oposición de género gramatical a
              través de la flexión.
          \ex Solo los sustantivos que designan entidades sexuadas expresan
              oposición de género a través de la flexión.
          \ex Solo los sustantivos que designan entidades sexuadas expresan
              oposición de género, ya sea mediante la flexión o por la
              concordancia con el determinante.
          \ex No todos los sustantivos expresan oposición de género gramatical.
        \end{exe}
  \item Con respecto a la realización de las oposiciones en género:
        \begin{exe}
          \ex La oposición de género gramatical se vincula siempre a que los
              hablantes de distintas variedades del español adscriben los
              sustantivos a distintos géneros.
          \ex La oposición de género gramatical se interpreta en todos los
              sustantivos como un contraste de sexo.
          \ex La oposición en género gramatical puede interpretarse como un
              contraste de rasgos descriptivos de la entidad designada por el
              tema o puede deberse a la variación lingüística.
          \ex La oposición de género gramatical se interpreta siempre como un
              contraste de rasgos descriptivos de la entidad designada por el
              tema.
        \end{exe}
  \item Con respecto a la realización de las oposiciones en género:
        \begin{exe}
          \ex La vocal \fonol{a} siempre forma la variante del sustantivo que
              concuerda con la variante femenina del determinante y la vocal
              \fonol{o} siempre forma la variante que concuerda con la variante
              masculina del determinante.
          \ex No hay una correspondencia estricta entre la vocal con la que
              termina un sustantivo y su género gramatical.
          \ex La vocal \fonol{a} siempre forma la variante del sustantivo que
              concuerda con la variante femenina del determinante y la variante
              masculina se puede expresar tanto con la vocal \fonol{o} como por
              la ausencia de vocal.
          \ex En los sustantivos que designan entidades sexuadas la vocal
              \fonol{a} siempre forma la variante del sustantivo que concuerda
              con la variante femenina del determinante y la vocal \fonol{o}
              siempre forma la variante que concuerda con la variante masculina
              del determinante.
        \end{exe}
\end{consignas}
\end{ejercicio}

\begin{ejercicio}\label{ej:04-verbos}
Considere el siguiente corpus con formas verbales de cada una de las tres
conjugaciones del español y luego resuelva las consignas de abajo:

\begin{tabular}{@{}lll@{}}
  \textbf{A} & \textbf{B} & \textbf{C} \\
  1. \textbf{llegaba} & 1. \textbf{comía} & 1. \textbf{aburría} \\
  2. llega     & 2. come     & 2. aburre \\
  3. llegaban  & 3. comían   & 3. aburrían \\
  4. llegó     & 4. comió    & 4. aburrió \\
  5. llegabas  & 5. comías   & 5. aburrías \\
  6. llegara   & 6. comiera  & 6. aburriera \\
\end{tabular}

\begin{consignas}
  \item Segmente en morfemas todas las palabras del corpus.
  \item Evalúe qué categorías gramaticales se oponen entre la forma 1 de cada
        columna y cada una de las demás de su misma columna.
  \item Haga la descripción morfológica de cada forma verbal, indique el
        contenido semántico de cada uno de los morfemas que la integran y
        complete el cuadro de abajo como se muestra en el modelo.
\end{consignas}

\begin{resuelto}[Ejemplos resueltos]
a) 1. A. \forma{lleg-a-ba}; B. \forma{com-í-a}; C. \forma{aburr-í-a}. 2. A.
\forma{lleg-a}; B. \forma{com-e}; C. \forma{aburr-e}. b) Contraste 1/2: ambas
formas se oponen en tiempo: 1) pretérito / 2) presente. c) 1.
\forma{llegaba/comía/aburría}: 1.ª/3.ª persona singular del pretérito imperfecto
del indicativo. 2. \forma{llega/come/aburre}: 3.ª persona singular del presente
de indicativo.
\end{resuelto}

{\footnotesize
\begin{tabularx}{\linewidth}{@{}l L L L l@{}}
  \toprule
  & Raíz & VT & TAM (amalgama de tiempo, modo y aspecto) & Conc. (número y persona) \\
  \midrule
  1. \forma{llegaba/comía/aburría} & lleg- / com- / aburr- & -a- / -í- & -ba / -a & ∅ \\
  2. \forma{llega/come/aburre}     & lleg- / com- / aburr- & -a / -e   &          & ∅ \\
  3. \forma{llegaban/comían/aburrían} & & & & \\
  4. \forma{llegó/comió/aburrió}   & & & & \\
  5. \forma{llegabas/comías/aburrías} & & & & \\
  6. \forma{llegara/comiera/aburriera} & & & & \\
  \bottomrule
\end{tabularx}}
\end{ejercicio}

\begin{recuadro}[Reflexión final]
La morfología es el componente de la gramática que, a partir de nuestra
competencia lingüística, involucra el conocimiento que nos permite segmentar las
palabras en unidades menores jerárquicamente relacionadas, establecer
vinculaciones formales entre las entradas que forman parte de nuestro léxico
mental y producir tanto elementos léxicos nuevos como las variantes gramaticales
que cada uno de ellos necesita exhibir en la combinatoria sintáctica. En este
sentido, aunque el nivel morfológico interactúa estrechamente con otros
componentes de la lengua, como la fonología, la sintaxis y el léxico, puede
pensarse que constituye un módulo diferenciado, delimitado por la naturaleza de
las unidades y los procesos que intervienen en él.
\end{recuadro}

% --- Secciones finales -------------------------------------------------
\lecturasrecomendadas

\solucionario

\begin{solucion}{\ref{ej:04-terraplanista}}
\textit{Parte A}. a) \forma{terraplanista}: \forma{terra-plan-ist-a};
\forma{negar}: \forma{neg-a-r}; \forma{más}: \forma{más}; \forma{falsedades}:
\forma{fals-e-dad-es}; \forma{gigantesca}: \forma{gigant-esc-a};
\forma{plenamente}: \forma{plen-a-mente}; \forma{explican}:
\forma{ex-plic-a-n}; \forma{entronización}: \forma{en-tron-iz-a-ción};
\forma{pseudocientíficos}: \forma{pseudo-cient-ífic-o-s}.

b) Se ofrecen ejemplos de respuestas posibles a modo ilustrativo.
\forma{terraplanista}: \forma{terra-}: \forma{terrá-que-o, terri-torio};
\forma{-plan-}: \forma{plan-o, plan-icie}; \forma{-ista}: \forma{tax-ista,
conform-ista}. \forma{negar}: \forma{neg-}: \forma{neg-ación, in-neg-able};
\forma{-a-}: \forma{am-a-r, explic-a-n}; \forma{-r}: \forma{teme-r, parti-r}.
\forma{falsedades}: \forma{fals-}: \forma{fals-o, fals-ificar};
\forma{-e-dad-}: \forma{sol-e-dad, ebri-e-dad}; \forma{-es}: \forma{lápic-es,
tabú-es}. \forma{gigantesca}: \forma{gigant-}: \forma{gigant-es,
a-gigant-a-r}; \forma{-esc-}: \forma{dant-esc-a, carnaval-esc-o}; \forma{-a}:
\forma{león-a, maestr-a}. \forma{entronización}: \forma{en-}:
\forma{en-amorar, en-casillaron}; \forma{-tron-}: \forma{des-tron-aron,
en-tron-izar}; \forma{-iz-}: \forma{real-iz-ó, esclav-iz-a-ron}; \forma{-a-}:
VT; \forma{-ción}: \forma{conten-ción, nega-ción}. \forma{pseudocientífico}:
\forma{pseudo-}: \forma{pseudo-Aristóteles, pseudo-efedrina}; \forma{-cient-}:
\forma{cient-ificismo, cient-ología}; \forma{-ífic-}: \forma{magn-ífic-o,
terror-ífic-as}; \forma{-o}: \forma{amarill-o, segur-o-s}.

c) \forma{con-venc-e-r}. Significado: \sig{mover con razones a alguien a hacer o
pensar algo}. Aporte semántico de cada morfema: \forma{con-}: no posee un
significado preciso en la palabra; tiene valor distintivo porque permite
distinguir el verbo \forma{convencer} de \forma{vencer}; \forma{-venc-}:
\sig{prevalecer sobre algo o alguien}; \forma{-e-}: VT; \forma{-r}: desinencia
de infinitivo. \forma{plen-a-mente}. Significado: \sig{de manera plena o
completa}. Aporte semántico de cada morfema: \forma{plen-}: \sig{lleno,
completo}; \forma{-a-}: no posee significado; es un resto etimológico del morfema
de género femenino que en estados previos de lengua marcaba la concordancia con
el sustantivo femenino \forma{mente}; \forma{-mente}: \sig{de manera
X}.\footnote{X remite al contenido del adjetivo: \textit{rápidamente} se
interpreta como \sig{de manera rápida} y \textit{alegremente}, \sig{de manera
alegre}.} \forma{ex-plic-a-n}. Significado: \sig{hacen comprender algo}. Aporte
semántico de cada morfema: \forma{ex-}: \sig{fuera}; \forma{-plic-}:
etimológicamente remite a la misma raíz que \forma{pleg-a-r}; la presencia del
prefijo \forma{ex-} produce la interpretación de conjunto \sig{desplegar hacia
afuera}, es decir, \sig{hacer conocer o comprender}; \forma{-a-}: VT;
\forma{-n}: desinencia de número y persona (amalgama, v.~§\ref{sec:04-vt}).

\textit{Parte B}. b) Opción correcta: 1. c) Opción correcta: 1.
\end{solucion}

\begin{solucion}{\ref{ej:04-alomorfos}}
Además del ejemplo resuelto, en el corpus se reconocen tres grupos de palabras
más:

1. a) En \forma{i-rrump-ido, rot-o, rupt-ura, i-rrup-ción, dis-rupt-iva,
romp-ió} se identifica el morfema \textsc{romp-}, que significa \sig{separar las
partes de un todo deshaciendo su unión}. Los alomorfos son \fonol{romp}
\textasciitilde{} \fonol{rot} \textasciitilde{} \fonol{rupt} \textasciitilde{}
\fonol{rup} \textasciitilde{} \fonol{rump}. b) El corpus no permite identificar
un factor que determine la aparición de cada variante. La causa de la alomorfía
es léxica.

2. a) En \forma{terra-planistas, terrá-queos, terr-ícola, en-tierr-an,
terr-itorio, a-terr-izar, Tierra} se identifica el morfema \textsc{terr-}, que
significa \sig{planeta que habitamos}. Los alomorfos son \fonol{ˈtjer}
\textasciitilde{} \fonol{ter}. b) La causa de la alomorfía es fonológica:
\fonol{ˈtjer} aparece en sílaba acentuada y \fonol{ter}, en sílaba no acentuada.

3. a) En \forma{terraplanista-s, cultural-es, terráqueo-s, dato-s,
falsedad-es, joven-es, maniquea-s, experto-s, (las) crisis} se identifica el
morfema \forma{-s}, que significa \sig{plural}. Los alomorfos son \fonol{s}
\textasciitilde{} \fonol{es} \textasciitilde{} ∅ (alomorfo cero). b) La causa de
la alomorfía es fonológica: \fonol{s} se adjunta a bases terminadas en vocal, con
excepción de las vocales cerradas acentuadas; y \fonol{es}, a bases terminadas en
vocal cerrada acentuada o en consonante, con excepción de las palabras graves y
esdrújulas terminadas en \fonol{s}, que quedan igual, por lo que en ellas se
manifiesta el alomorfo cero, sin contenido fonológico.
\end{solucion}

\begin{solucion}{\ref{ej:04-clases}}
\textit{Parte A}. a)

{\footnotesize
\begin{tabular}{@{}l*{5}{c}@{}}
  \toprule
  Palabra & \makecell{Pref.\\flexivo} & \makecell{Pref.\\derivativo} & Raíz & \makecell{Suf.\\derivativo} & \makecell{Suf.\\flexivo} \\
  \midrule
  \forma{terraplanista}     & & & terra- -plan- & -ista & \\
  \forma{negar}             & & & neg- & -a- & -r \\
  \forma{más}               & & & más & & \\
  \forma{achatada}          & & a- & -chat- & -a- -d- & -a \\
  \forma{gigantesca}        & & & gigant- & -esc- & -a \\
  \forma{plenamente}        & & & plen- & -a- -mente & \\
  \forma{representativos}   & & re- & -present- & -a- -t- -iv- & -o- -s \\
  \forma{explican}          & & ex- & -plic- & -a- & -n \\
  \forma{entronización}     & & en- & -tron- & -iz- -a- -ción & \\
  \forma{pseudocientíficos} & & pseudo- & -cient- & -ífic- & -o- -s \\
  \bottomrule
\end{tabular}}

b) El único morfema libre en el corpus es \forma{más}. El resto de los formantes
son ligados. \textit{Parte B}. b) Opción correcta: 3.
\end{solucion}

\begin{solucion}{\ref{ej:04-tema}}
\textit{Parte A}. a) \forma{culturales}: \forma{cult-ur-al-es} = \sig{relativos a
la cultura}; \forma{bromear}: \forma{brom-e-a-r} = \sig{hacer bromas};
\forma{desprecio}: \forma{des-precio} = \sig{falta de aprecio}.

b) \forma{chato}\textsubscript{A} > \forma{achata(r)}\textsubscript{V} >
\forma{achatada}\textsubscript{A};\footnote{\textit{Achatada} es el participio
del V \textit{achatar} que, por un proceso de conversión, puede funcionar
también como adjetivo (v.~§\ref{sec:04-conversion}).} A > V > Part/A. Adjetivo:
\forma{La Tierra es una esfera achatada}. \forma{chato}\textsubscript{A} >
\forma{achata(r)}\textsubscript{V} > \forma{achatada}\textsubscript{V part.}; A
> V > V\textsubscript{part.}. Participio: \forma{La chapa fue achatada con una
maza por el herrero}. \forma{ciencia}\textsubscript{N} >
\forma{cien(cia)tífico}\textsubscript{A} >
\forma{pseudocientífico}\textsubscript{A}; N > A > A.
\forma{culto}\textsubscript{A} > \forma{cult(o)ura}\textsubscript{N} >
\forma{cultur(a)ales}\textsubscript{A}; A > N > A.
\forma{broma}\textsubscript{N} > \forma{brom(a)ea(r)}\textsubscript{V}; N > V.
\forma{precio}\textsubscript{N} > \forma{desprecio}\textsubscript{N}; N >
N\textsubscript{neg}.

c) DE:
[[a-[chat]\textsubscript{A}-a]\textsubscript{V}-d-a]\textsubscript{A} (como
adjetivo); [a-[chat]\textsubscript{A}-a-d-a]\textsubscript{V} (como participio
verbal);
[pseudo-[[cient]\textsubscript{N}-ífic-o]\textsubscript{A}]\textsubscript{A};
[[[cult]\textsubscript{A}-ur]\textsubscript{N}-al]\textsubscript{A};
[[brom]\textsubscript{N}-e-a-r]\textsubscript{V};
[des-[precio]\textsubscript{N}]\textsubscript{N}.

\textit{Parte B}. b) Opción correcta: 3.
\end{solucion}

\begin{solucion}{\ref{ej:04-homonimia}}
2. \forma{inutilizable}: a. \sig{que se puede inutilizar}; b. \sig{que no se
puede utilizar}. 3. \forma{desarmar}: a. \sig{revertir el proceso de armar (un
discurso)}; b. \sig{quitarle el arma a alguien}.

2. a. \forma{útil}\textsubscript{A} > \forma{inútil}\textsubscript{A} >
\forma{inutilizar}\textsubscript{V} > \forma{inutilizable}\textsubscript{A};
[[[in-[útil]\textsubscript{A}]\textsubscript{A}-iza]\textsubscript{V}-ble]\textsubscript{A};
A > A > V > A. b. \forma{útil}\textsubscript{A} >
\forma{utilizar}\textsubscript{V} > \forma{utilizable}\textsubscript{A} >
\forma{inutilizable}\textsubscript{A};
[in-[[[útil]\textsubscript{A}-iza]\textsubscript{V}-ble]\textsubscript{A}]\textsubscript{A};
A > V > A > A.

3. a. \forma{armar}\textsubscript{V} > \forma{desarmar}\textsubscript{V};
[des-[arm-a-r]\textsubscript{V}]\textsubscript{V}; V > V. b.
\forma{arma}\textsubscript{N} > \forma{armar}\textsubscript{V} >
\forma{desarmar}\textsubscript{V}; [des-[arm]\textsubscript{N}-a-r]\textsubscript{V};
N > V > V.
\end{solucion}

\begin{solucion}{\ref{ej:04-ardo}}
\textit{Parte A}. b) \forma{-ard-}: aunque el \textit{DLE} sostiene que este
sufijo forma sustantivos y adjetivos con valor aumentativo o despectivo, en el
uso juvenil actual su empleo asigna \sig{valoración positiva por parte del
hablante}. c) Ejemplos posibles sugeridos: \forma{epic-ard-o, amig-ard-o,
cervez-ard-a, vin-ard-o, mes-ard-a, min-ard-a, tip-ard-o, te-s-ard-o}.

\textit{Parte B}. a) 2; b) 2; c) 1; d) 4. El aporte semántico de \forma{-ard-} se
asocia a significados expresivos vinculados con valoraciones del hablante.

\textit{Parte C}. Las características evaluadas en los puntos anteriores
permiten incluir el sufijo \forma{-ardo} entre los que pertenecen a la
morfología apreciativa.
\end{solucion}

\begin{solucion}{\ref{ej:04-derivacion}}
\textit{Parte A}. Grupo 2. a) \forma{falsedades, habilidad, deslealtad,
cotidianidad, bondades, maldad, disconformidad}. Proceso común: derivación con
el sufijo \forma{-dad}. b) Categoría de la base: A. Categoría resultante: N. c)
Valor semántico: \sig{propiedad de ser X}. d) Descripción fonológica:
\fonol{ˈdad}, \fonol{ˈtad}. Las vocales \fonol{e} e \fonol{i} en \fonol{ˈedad} e
\fonol{ˈidad} pueden considerarse interfijos cuya aparición se explica por
razones fonológicas.

Grupo 3. a) \forma{chateaba, bromear, desposteó, googlearon}. Proceso común:
derivación con el sufijo \forma{-e(a)-}. b) Categoría de la base: N. Categoría
resultante: V. c) Valor semántico: \sig{hacer X}. d) Descripción fonológica:
\fonol{eˈa}.

Grupo 4. a) \forma{deslealtad, desinformó, desposteó, desagradables, desamor,
disconformidad}. Proceso común: derivación con el prefijo \forma{des-}. b)
Categorías de la base: A\textsubscript{x}, N\textsubscript{y},
V\textsubscript{z}. Categoría resultante: A\textsubscript{x},
N\textsubscript{y}, V\textsubscript{z} (los índices señalan que las bases no
cambian su categoría). c) Valor semántico: \sig{inversión / negación / privación
/ contrariedad respecto de X}. d) Descripción fonológica: \fonol{ˈdes},
\fonol{ˈdis}.

\textit{Parte B}. Opción correcta: 2.
\end{solucion}

\begin{solucion}{\ref{ej:04-composicion}}
Grupo 2. a) \forma{sacacorchos, lavarropas, guardapolvos, pisapapeles,
pelapapas, cazarrecompensas, matabichos, cometierra, chupamedias,
limpiaparabrisas}. b) \forma{saca-corchos; lava-ropas; guarda-polvos;
pisa-papeles; pela-papas; caza-recompensas; mata-bichos; come-tierra;
chupa-medias; limpia-parabrisas}. c)
[[saca]\textsubscript{V}[corchos]\textsubscript{N}]\textsubscript{N}. d) V + N >
N.

Grupo 3. a) \forma{sordomudo, socioeconómico, blanquiceleste, dentoalveolar,
italoargentino, agridulce}. b) \forma{sordo-mudo; socio-económico;
blanqu-i-celeste; dent-o-alveolar; ital-o-argentino; agri-dulce}. c)
[[sordo]\textsubscript{A}[mudo]\textsubscript{A}]\textsubscript{A}. d) A + A > A.
\end{solucion}

\begin{solucion}{\ref{ej:04-filos}}
\textit{Parte A}. \forma{filosofía}:
[[filo]\textsubscript{TC}[sofía]\textsubscript{TC}]\textsubscript{N}, TC + TC >
N; \forma{melófila}: [[meló]\textsubscript{TC}[fil]\textsubscript{TC}a]\textsubscript{A},
TC + TC > A; \forma{filonazi}:
[[filo]\textsubscript{TC}[nazi]\textsubscript{N}]\textsubscript{A}, TC + N > A;
\forma{hidrófilo}:
[[hidró]\textsubscript{TC}[fil]\textsubscript{TC}o]\textsubscript{A}, TC + TC >
A; \forma{americanófila}:
[[american]\textsubscript{N}ó[fil]\textsubscript{TC}a]\textsubscript{A}, N + TC >
A; \forma{filología}:
[[filo]\textsubscript{TC}[logía]\textsubscript{TC}]\textsubscript{N}, TC + TC >
N; \forma{tecnófilo}:
[[tecno]\textsubscript{TC}[fil]\textsubscript{TC}o]\textsubscript{A}, TC + TC >
A; \forma{filántropo}:
[[fil]\textsubscript{TC}[antrop]\textsubscript{TC}o]\textsubscript{N}, TC + TC >
N.

\textit{Parte B}. Opciones correctas: a) 3; b) 1.
\end{solucion}

\begin{solucion}{\ref{ej:04-parasintesis}}
\pendiente{El solucionario del original no corresponde a las consignas de este
ejercicio (analiza afijos y no bases parasintéticas); revisar}
Se ofrecen ejemplos de respuestas posibles a lo pedido a modo ilustrativo.
\forma{terra-plan-ista}. \forma{-ista}: a) \forma{tech-ista, social-ista}; b)
categoría de la base: N/A; c) categoría del resultado: N; N/A + \forma{-ista} >
N; d) valor semántico: \sig{que profesa / practica / cree en X}.
\forma{sub-terr-áne-a}. \forma{sub-}: a) \forma{sub-suelo, sub-valorar}; b)
categoría de la base: N/A/V; c) categoría del resultado: N/A/V;
X\textsubscript{i} + \forma{sub-} > X\textsubscript{i}; d) valor semántico:
\sig{debajo de, menos que X}. \forma{-áne-}: a) \forma{moment-áne-o,
cut-áne-a}; b) categoría de la base: N; c) categoría del resultado: A; N +
\forma{-áne-} > A; d) valor semántico: \sig{relacionado con / relativo a X}.
\forma{mil-eur-ista-s}. \forma{-ista} (v. \forma{terraplanista}).
\forma{sal-piment-a-r}. \forma{-a-}: a) \forma{lider-a-r, activ-a-r}; b)
categoría de la base: N/A; c) categoría del resultado: V; N/A + \forma{-a-} > V;
d) valor semántico: \sig{usar X}. \forma{in-fum-a-ble}. \forma{in-}: a)
\forma{in-moral, in-útil}; b) categoría de la base: A; c) categoría del
resultado: A; \forma{in-} + A > A; d) valor semántico: \sig{no X}.
\forma{-ble}: a) \forma{cont-a-ble, tem-i-ble-s}; b) categoría de la base: V; c)
categoría del resultado: A; V + \forma{-ble} > A; d) valor semántico: \sig{que
se puede X}.
\end{solucion}

\begin{solucion}{\ref{ej:04-conversion}}
En todos los casos las formas comparten la raíz y lo que varía es el tema.
2. \forma{com-i-d-a}. a) \forma{comida}\textsubscript{a}: tema \forma{comida-};
\forma{comida}\textsubscript{b}: tema \forma{comi-}. b)
\forma{comida}\textsubscript{a} es sustantivo porque va precedido del
determinante \forma{la}; \forma{comida}\textsubscript{b} es una forma verbal de
participio del verbo \forma{comer} que integra una perífrasis de voz pasiva.
3. \forma{curs-a-d-a}. a) \forma{cursada}\textsubscript{a}: tema
\forma{cursada-}; \forma{cursada}\textsubscript{b}: tema \forma{cursa-}. b)
\forma{cursada}\textsubscript{a} funciona como sustantivo porque va precedido
del determinante \forma{la}; \forma{cursada}\textsubscript{b} es la forma verbal
de participio del verbo \forma{cursar} empleada como adjetivo (por conversión)
que modifica al sustantivo \forma{notificación}.\pendiente{Revisar temas: el
original los invierte en a y b}
4. \forma{priv-a-d-o}. a) \forma{privado}\textsubscript{a}: tema
\forma{privad-}; \forma{privado}\textsubscript{b}: tema \forma{priva-}. b)
\forma{privado}\textsubscript{a} funciona como un adjetivo que modifica a
\forma{empleo}; \forma{privado}\textsubscript{b} es una forma verbal de
participio del verbo \forma{privar} que integra una perífrasis de voz pasiva.
5. \forma{doctor-a-ndo}. a) \forma{doctorando}\textsubscript{a}: tema
\forma{doctora-}; \forma{doctorando}\textsubscript{b}: tema
\forma{doctorando-}. b) \forma{doctorando}\textsubscript{a} es la forma verbal
de gerundio del verbo \forma{doctorarse} que integra una perífrasis verbal
progresiva de \forma{estar} + gerundio; \forma{doctorando}\textsubscript{b}
funciona como un sustantivo porque va precedido por el determinante \forma{los}.
6. \forma{cumbia}. a) \forma{cumbia}\textsubscript{a}: tema \forma{cumbia-};
\forma{cumbia}\textsubscript{b}: tema \forma{cumbia-}. b)
\forma{cumbia}\textsubscript{a} funciona como un sustantivo precedido por el
determinante \forma{esta}; \forma{cumbia}\textsubscript{b} ejemplifica un proceso
de conversión similar al del inglés, en el que los sustantivos pueden dar origen
a verbos de significado relacionado.
\end{solucion}

\begin{solucion}{\ref{ej:04-regresiva}}
2. \forma{espectar}. a) Base: \forma{espectador-}. Categoría: N.\pendiente{El
original dice «A»} Resultado: \forma{espect(a)-}. Categoría: V. Segmento
afectado: \forma{-dor}. b) \forma{espectar} < \forma{espectador}: N --
\forma{-dor} > V. c) Significado: \sig{participar de un juego como espectador}.
3. \forma{alfabeto}. a) Base: \forma{analfabet-}. Categoría: A. Resultado:
\forma{alfabet-}. Categoría: A. Segmento afectado: \forma{an-}. b)
\forma{alfabeto} < \forma{analfabeto}: A -- \forma{an-} > A. c) Significado:
\sig{que no es analfabeto}. 4. \forma{antitranspirar}. a) Base:
\forma{antitranspirante-}. Categoría: A. Resultado: \forma{antitranspir(a)-}.
Categoría: V. Segmento afectado: \forma{-nte}. b) \forma{antitranspirar} <
\forma{antitranspirante}: A -- \forma{-nte} > V. c) Significado: \sig{acción
contraria a transpirar}.
\end{solucion}

\begin{solucion}{\ref{ej:04-acortamiento}}
Grupo 2: acortamientos por apócope. a) \forma{profe, facu, auto, moto, compu,
taxi}. b) \forma{profe}\textsubscript{N} < \forma{profesor/a}\textsubscript{N};
\forma{facu}\textsubscript{N} < \forma{facultad}\textsubscript{N};
\forma{auto}\textsubscript{N} < \forma{automóvil}\textsubscript{N};
\forma{moto}\textsubscript{N} < \forma{motocicleta}\textsubscript{N};
\forma{compu}\textsubscript{N} < \forma{computadora}\textsubscript{N};
\forma{taxi}\textsubscript{N} < \forma{taxímetro}\textsubscript{N}. c) Se
conservan las primeras dos sílabas de la base. El resultado, al igual que en el
caso del grupo 1, tiene acentuación grave. d) El acortamiento por apócope
tampoco produce cambio alguno en el significado de la base.
\end{solucion}

\begin{solucion}{\ref{ej:04-cruces}}
2. \forma{Conmebol}. a) Bases: \forma{confederación + sudamericana + (de)
fútbol}. b) Categoría de las bases: N + A + N. Categoría del resultado: N. c)
Proceso de selección: primera sílaba de la primera base, tercera sílaba de la
segunda base, última sílaba de la última base: [con(federación) +
(suda)me(ricana) + (fút)bol]. 3. \forma{Mercosur}. a) Bases: \forma{mercado +
común + (del) sur}. b) Categoría de las bases: N + A + N. Categoría del
resultado: N. c) Proceso de selección: primera sílaba de cada base; la última de
las bases es monosilábica y se integra completa: [mer(cado) + co(mún) + (del)
sur]. 4. \forma{vitalitud}. a) Bases: \forma{vitalidad + actitud}. b) Categoría
de las bases: N + N. Categoría del resultado: N. c) Proceso de selección:
primeros seis segmentos sonoros de la primera base y cuatro segmentos finales de
la segunda, aprovechando el fonema vocálico común: [vital(idad) + (acti)tud]. 5.
\forma{deshamburguesar}. a) Bases: \forma{desaburguesar + hamburguesa}. b)
Categoría de las bases: V + N. Categoría del resultado: V. c) Proceso de
selección: primeros cuatro segmentos sonoros de la primera base y ocho segmentos
finales de la segunda, aprovechando la coincidencia en los últimos siete
fonemas: [desa(burguesar) + (hamburgues)a]. En la publicidad, el verbo
\forma{deshamburguesar} aparece en la forma imperativa de la 2.ª persona
singular, por lo que no se expresan morfemas flexivos, facilitando la
coincidencia entre ambas bases. En otras formas, \forma{deshamburguesar} debe
incorporar la morfología flexiva correspondiente a la 1.ª conjugación verbal,
ausente en el sustantivo \forma{hamburguesa}. 6. \forma{artivismo}. a) Bases:
\forma{arte + activismo}. b) Categoría de las bases: N + N. Categoría del
resultado: N. c) Proceso de selección: primeros tres segmentos sonoros de la
primera base y siete segmentos finales de la segunda, aprovechando el fonema
consonántico común: [art(e) + (ac)tivismo]. 7. \forma{juernes}. a) Bases:
\forma{jueves + viernes}. b) Categoría de las bases: N + N. Categoría del
resultado: N. c) Proceso de selección: primeros tres segmentos sonoros de la
primera base y cuatro segmentos finales de la segunda: [jue(ves) + (vie)rnes].
\end{solucion}

\begin{solucion}{\ref{ej:04-acronimos}}
En todos los casos el acrónimo equivale referencialmente al sintagma que
abrevia. 2. \forma{RAE}. a) Real\textsubscript{A} Academia\textsubscript{N}
Española\textsubscript{A}. b) Categoría de la construcción: N. c) No posee marcas
morfológicas adicionales. 3. \forma{ovnis}. a) objetos\textsubscript{N}
voladores\textsubscript{A} no\textsubscript{IP}\footnote{La sigla IP corresponde
a \textit{índice de polaridad}, en este caso, negativa. El concepto de
polaridad se desarrolla en el \capref{14}.} identificados\textsubscript{A}. b)
Categoría del resultado: N. c) El segmento final corresponde con la
pluralización, que se manifiesta por medio del alomorfo correspondiente a la
terminación del tema (\forma{-s}). 4. \forma{FIFA / AFA}. a)
Federación\textsubscript{N} Internacional\textsubscript{A} (de)
Fútbol\textsubscript{N} Asociado\textsubscript{A} / Asociación\textsubscript{N}
(del) Fútbol\textsubscript{N} Argentino\textsubscript{A}. b) Categoría de los
resultados: N. c) No poseen marcas morfológicas adicionales. 5. \forma{UNA}. a)
Universidad\textsubscript{N} Nacional\textsubscript{A} (de las)
Artes\textsubscript{N}. b) Categoría del resultado: N. c) No posee marcas
morfológicas adicionales.
\end{solucion}

\begin{solucion}{\ref{ej:04-siempre}}
a) Categoría de las bases: A. Categoría del resultado: A. b) El segmento
\forma{in-} con que comienzan todas estas palabras se reinterpreta como si fuera
el prefijo negativo \forma{in-}. c) En \forma{imberbe}, \forma{im-} representa
al prefijo negativo, ya que la palabra significa \sig{que no tiene barba}; no
obstante, \forma{berbe} constituye un caso de retroformación, dado que esta
palabra no existe en español y se formaría a partir de \forma{imberbe}.
\end{solucion}

\begin{solucion}{\ref{ej:04-reanalisis}}
Los tres son casos de reanálisis. 2. \forma{dictablanda}. Se trata de la
reinterpretación de la desinencia \forma{-dura} de \forma{dictadura}, formante de
sustantivos abstractos (\forma{gordura, dulzura}), como si se tratara del
adjetivo \forma{duro/a}, lo que permite la formación de un supuesto opuesto:
\forma{dictablanda}. El término se usó para hacer referencia al gobierno del
general Berenguer en España, ocurrido tras la dictadura de Primo de Rivera. Como
el gobierno de Berenguer se inició con algunas medidas liberalizadoras, se lo
denominó \forma{dictablanda}. 3. \forma{precuela}: versión española del inglés
\forma{prequel}, formado a partir de \forma{sequel} \sig{secuela}, derivado del
verbo latino \forma{sequor} (\sig{seguir}), sobre una base inexistente
\juicio{*}{\forma{-quel}} a la que se le añade el prefijo \forma{pre-} para
hacer referencia a una «obra literaria o cinematográfica que cuenta hechos que
preceden a los de otra obra ya existente» (\textit{DLE}, en línea).
\end{solucion}

\begin{solucion}{\ref{ej:04-genero}}
\textit{Parte A}. 2. \forma{el sartén / la sartén}. 3. \forma{la gat-a / el
gat-o}. 4. \forma{la cebra} / ∅. 5. \forma{el naranj-o / la naranj-a}. 6.
\forma{la cuchill-a / el cuchill-o}. 7. \forma{la araña} / ∅. 8. \forma{la
mano} / ∅. 9. \forma{el clima} / ∅. 10. \forma{el alcalde / la alcald-esa}. 11.
\forma{el amante / la amante}. 12. \forma{el tilde / la tilde}. 13. \forma{la
bols-a / el bols-o}. 14. \forma{el calor / la calor}.

\textit{Parte B}. a) Opción correcta: 4. Sustantivos como \forma{cebra},
\forma{araña}, \forma{mano} y \forma{clima} poseen género inherente. A los
sustantivos como \forma{cebra} y \forma{araña}, que abarcan indistintamente
individuos machos o hembras de la clase de animal denotada, se los conoce como
epicenos. b) Opción correcta: 3. Además de la oposición en términos de
macho/hembra para sustantivos que designan entidades sexuadas (como
\forma{secretario/secretaria}), la categoría también abarca otros contrastes de
significado, como árbol/fruto (\forma{naranjo/naranja}) o tamaño
(\forma{cuchillo/cuchilla, bolso/bolsa}). Asimismo, la selección de género
inherente en algunos nombres puede estar sujeta a variación socioestilística
(\forma{el/la sartén, el/la tilde}). c) Opción correcta: 2.
\end{solucion}

\begin{solucion}{\ref{ej:04-verbos}}
a) 3. A. \forma{lleg-a-ba-n}; B. \forma{com-í-a-n}; C. \forma{aburr-í-a-n}. 4.
A. \forma{lleg-ó}; B. \forma{com-i-ó}; C. \forma{aburr-i-ó}. 5. A.
\forma{lleg-a-ba-s}; B. \forma{com-í-a-s}; C. \forma{aburr-í-a-s}. 6. A.
\forma{lleg-a-ra}; B. \forma{com-ie-ra}; C. \forma{aburr-ie-ra}.

b) Contraste 1/3: oposición de número: 1) singular / 3) plural. Contraste 1/4:
oposición de aspecto: 1) imperfecto / 4) perfecto. Contraste 1/5: oposición de
persona: 1) 1.ª/3.ª / 5) 2.ª. Contraste 1/6: oposición de modo: 1) indicativo /
6) subjuntivo.

c) 3. 3.ª persona plural del pretérito imperfecto del indicativo. 4. 3.ª persona
singular del pretérito perfecto simple del indicativo. 5. 2.ª persona singular
del pretérito imperfecto del indicativo. 6. 1.ª/3.ª persona singular del
pretérito imperfecto del subjuntivo.

{\footnotesize
\begin{tabularx}{\linewidth}{@{}l L L L l@{}}
  \toprule
  & Raíz & VT & TAM & Conc. \\
  \midrule
  1. \forma{llegaba/comía/aburría}     & lleg- / com- / aburr- & -a- / -í- & -ba / -a & ∅ \\
  2. \forma{llega/come/aburre}         & lleg- / com- / aburr- & -a / -e   &          & ∅ \\
  3. \forma{llegaban/comían/aburrían}  & lleg- / com- / aburr- & -a- / -í- & -ba / -a & -n \\
  4. \forma{llegó/comió/aburrió}       & lleg- / com- / aburr- & ∅ / -i-   & -ó       & ∅ \\
  5. \forma{llegabas/comías/aburrías}  & lleg- / com- / aburr- & -a- / -í- & -ba / -a & -s \\
  6. \forma{llegara/comiera/aburriera} & lleg- / com- / aburr- & -a- / -ie- & -ra     & ∅ \\
  \bottomrule
\end{tabularx}}
\end{solucion}
